\begingroup
\fontsize{12}{20.7}\selectfont
\raggedright
\begin{list}{}{\setlength{\leftmargin}{10mm}\setlength{\labelwidth}{8mm}\setlength{\labelsep}{2mm}\setlength{\itemsep}{0pt}\setlength{\parsep}{0pt}\setlength{\topsep}{0pt}}
\item[1.] A canonical excitatory–inhibitory motif embedded in the connectome underlies rhythmic locomotion in Drosophila larvae : препринт / T. Seki, H. Kohsaka, M. F. Zwart, A. Nose. — bioRxiv. — 2026. — DOI 10.64898/2026.08.31.748252. — URL: \url{https://www.biorxiv.org/content/10.64898/2026.08.31.748252v1.full} (дата обращения: 28.09.2026). — Текст : электронный.
\item[2.] A Clinical Feasibility Study of Spinal Evoked Compound Action Potential Estimation Methods / K. Chakravarthy, J. FitzGerald, A. Will [и др.] // Neuromodulation: Technology at the Neural Interface. — 2022. — Vol. 25, no. 1. — P. 75–84. — DOI 10.1111/ner.13510. — URL: \url{https://pubmed.ncbi.nlm.nih.gov/35041590/} (дата обращения: 01.10.2026). — Текст : электронный.
\item[3.] A Connectome of the Male Drosophila Ventral Nerve Cord : препринт / S. Y. Takemura, K. J. Hayworth, G. B. Huang [и др.]. — eLife, reviewed preprint, версия 1. — 2024. — Vol. 13. — Art. RP97769. — DOI 10.7554/eLife.97769.1. — URL: \url{https://elifesciences.org/reviewed-preprints/97769} (дата обращения: 08.10.2026). — Текст : электронный.
\item[4.] A CrossMod-Transformer deep learning framework for multi-modal pain detection through EDA and ECG fusion / J. Farmani, G. Bargshady, S. Gkikas [и др.] // Scientific Reports. — 2025. — Vol. 15, no. 1. — Art. 29467. — DOI 10.1038/s41598-025-14238-y. — URL: \url{https://www.nature.com/articles/s41598-025-14238-y} (дата обращения: 08.10.2026). — Текст : электронный.
\item[5.] A descending inhibitory mechanism of nociception mediated by an evolutionarily conserved neuropeptide system in Drosophila / I. Oikawa, S. Kondo, K. Hashimoto [и др.] // eLife. — 2023. — Vol. 12. — Art. RP85760. — DOI 10.7554/eLife.85760. — URL: \url{https://www.ebi.ac.uk/europepmc/webservices/rest/PMC10328526/fullTextXML} (дата обращения: 27.09.2026). — Текст : электронный.
\item[6.] A Drosophila computational brain model reveals sensorimotor processing / P. K. Shiu, G. R. Sterne, N. Spiller [и др.] // Nature. — 2024. — Vol. 634, no. 8032. — P. 210–219. — DOI 10.1038/s41586-024-07763-9. — URL: \url{https://www.nature.com/articles/s41586-024-07763-9} (дата обращения: 25.09.2026). — Текст : электронный.
\item[7.] A generalized activating function for rapid and accurate neural response prediction in spinal cord stimulation / T. H. Newton, J. García Ordóñez, A. Alashqar [и др.] // Communications Biology. — 2026. — DOI 10.1038/s42003-026-10849-x. — URL: \url{https://doi.org/10.1038/s42003-026-10849-x} (дата обращения: 30.09.2026). — Текст : электронный.
\item[8.] A lightweight data-driven spiking neuronal network model of Drosophila olfactory nervous system with dedicated hardware support / T. Nanami, D. Yamada, M. Someya [и др.] // Frontiers in Neuroscience. — 2024. — Vol. 18. — Art. 1384336. — DOI 10.3389/fnins.2024.1384336. — URL: \url{https://pubmed.ncbi.nlm.nih.gov/38994271/} (дата обращения: 28.09.2026). — Текст : электронный.
\item[9.] A modular, high-bandwidth, bidirectional implantable device for neural interrogation / R. Darie, S. R. Parker, J. S. Calvert [и др.] // Journal of Neural Engineering. — 2026. — Vol. 23, no. 4. — Art. 046020. — DOI 10.1088/1741-2552/ae7f52. — URL: \url{https://iopscience.iop.org/article/10.1088/1741-2552/ae7f52} (дата обращения: 30.09.2026). — Текст : электронный.
\item[10.] A Multi-Modal Multi-Expert Framework for Pain Assessment in Postoperative Children / Z. Liang, H. Luo, X. Chen [и др.] // IEEE Transactions on Affective Computing. — 2025. — Vol. 16, no. 4. — P. 2828–2841. — DOI 10.1109/TAFFC.2025.3567307. — URL: \url{https://so-link.github.io/assets/pdf/A\_Multi-Modal\_Multi-Expert\_Framework\_for\_Pain\_Assessment\_in\_Postoperative\_Children.pdf} (дата обращения: 28.09.2026). — Текст : электронный.
\item[11.] A multilevel multimodal circuit enhances action selection in Drosophila / T. Ohyama, C. M. Schneider-Mizell, R. D. Fetter [и др.] // Nature. — 2015. — Vol. 520, no. 7549. — P. 633–639. — DOI 10.1038/nature14297. — URL: \url{https://www.nature.com/articles/nature14297} (дата обращения: 25.09.2026). — Текст : электронный.
\item[12.] A neural correlate of individual odor preference in Drosophila / M. A. Churgin, D. O. Lavrentovich, M. A. Y. Smith [и др.] // eLife. — 2025. — Vol. 12. — Art. RP90511. — DOI 10.7554/eLife.90511. — URL: \url{https://elifesciences.org/articles/90511} (дата обращения: 28.09.2026). — Текст : электронный.
\item[13.] A Neuromodulable Current-Mode Silicon Neuron for Robust and Adaptive Neuromorphic Systems : препринт / L. Mendolia, C. Wen, E. Chicca [и др.]. — arXiv. — 2025. — arXiv:2512.01133. — URL: \url{https://arxiv.org/abs/2512.01133} (дата обращения: 08.10.2026). — Текст : электронный.
\item[14.] A neuromorphic model of olfactory processing and sparse coding in the Drosophila larva brain / A. M. Jürgensen, A. Khalili, E. Chicca [и др.] // Neuromorphic Computing and Engineering. — 2021. — Vol. 1, no. 2. — Art. 024008. — DOI 10.1088/2634-4386/ac3ba6. — URL: \url{https://research.rug.nl/en/publications/a-neuromorphic-model-of-olfactory-processing-and-sparse-coding-in/} (дата обращения: 28.09.2026). — Текст : электронный.
\item[15.] A neuropeptidergic circuit gates selective escape behavior of Drosophila larvae / B. N. Imambocus, F. Zhou, A. Formozov [и др.] // Current Biology. — 2022. — Vol. 32, no. 1. — P. 149–163.e8. — DOI 10.1016/j.cub.2021.10.069. — URL: \url{https://nms.unl.pt/Portals/0/adam/NMS\%20Generic\%20Content/cmXys54yHUWJBiTqEQLxqA/HighlightedPublications/Imambocus\%20et\%20al\%20-\%20Curr\%20Biol\%202021.pdf} (дата обращения: 27.09.2026). — Текст : электронный.
\item[16.] A novel thermosensitive escape behavior in Drosophila larvae / M. Oswald, B. Rymarczyk, A. Chatters, S. Sweeney // Fly. — 2011. — Vol. 5, no. 4. — P. 304–306. — DOI 10.4161/fly.5.4.17810. — URL: \url{https://pmc.ncbi.nlm.nih.gov/articles/PMC3266071/} (дата обращения: 26.09.2026). — Текст : электронный.
\item[17.] A positive feedback loop between sensory and octopaminergic neurons underlies nociceptive plasticity in Drosophila larvae / J. C. Boivin, Y. Q. Zhao, J. Zhu [и др.] // PLOS Genetics. — 2026. — Vol. 22, no. 4. — Art. e1012122. — DOI 10.1371/journal.pgen.1012122. — URL: \url{https://pubmed.ncbi.nlm.nih.gov/42048383/} (дата обращения: 08.10.2026). — Текст : электронный.
\item[18.] A predictive corticospinal model for pain perception / X. M. Lin, X. S. Zhang, H. Zhou [и др.] // Cell Reports Medicine. — 2026. — Vol. 7, no. 7. — Art. 102793. — DOI 10.1016/j.xcrm.2026.102793. — URL: \url{https://doi.org/10.1016/j.xcrm.2026.102793} (дата обращения: 08.10.2026). — Текст : электронный.
\item[19.] A programmable and self-adaptive ultrasonic wireless implant for personalized chronic pain management / Y. Zeng, C. Gong, G. Lu [и др.] // Nature Electronics. — 2025. — Vol. 8, no. 5. — P. 437–449. — DOI 10.1038/s41928-025-01374-6. — URL: \url{https://pmc.ncbi.nlm.nih.gov/articles/PMC13095316/} (дата обращения: 24.09.2026). — Текст : электронный.
\item[20.] A Single-Cell Level and Connectome-Derived Computational Model of the Drosophila Brain / Y. C. Huang, C. T. Wang, T. S. Su [и др.] // Frontiers in Neuroinformatics. — 2019. — Vol. 12. — Art. 99. — DOI 10.3389/fninf.2018.00099. — URL: \url{https://www.frontiersin.org/journals/neuroinformatics/articles/10.3389/fninf.2018.00099/full} (дата обращения: 28.09.2026). — Текст : электронный.
\item[21.] A Systematic Review of Multimodal Signal Fusion for Acute Pain Assessment Systems / M. U. Khan, G. Chetty, R. Goecke, R. Fernandez-Rojas // ACM Computing Surveys. — 2026. — Vol. 58, no. 2. — P. 1–35. — DOI 10.1145/3737281. — URL: \url{https://researchsystem.canberra.edu.au/ws/portalfiles/portal/118630433/3737281.pdf} (дата обращения: 22.09.2026). — Текст : электронный.
\item[22.] A wearable pain relief device and remote monitoring system therefor : патент CN120478838B, Китай. — 2026. — URL: \url{http://epub.cnipa.gov.cn/patent/CN120478838B} (дата обращения: 24.09.2026). — Текст : электронный.
\item[23.] Abdel Deen, O. M. T. A Multimodal Deep Learning Approach to Intraoperative Nociception Monitoring: Integrating Electroencephalogram, Photoplethysmography, and Electrocardiogram / O. M. T. Abdel Deen, S. Z. Fan, J. S. Shieh // Sensors. — 2025. — Vol. 25, no. 4. — Art. 1150. — DOI 10.3390/s25041150. — URL: \url{https://pmc.ncbi.nlm.nih.gov/articles/PMC11859842/} (дата обращения: 22.09.2026). — Текст : электронный.
\item[24.] Absence of paresthesia during high-rate spinal cord stimulation reveals importance of synchrony for sensations evoked by electrical stimulation / B. Sagalajev, T. Zhang, N. Abdollahi [и др.] // Neuron. — 2024. — Vol. 112, no. 3. — P. 404–420.e6. — DOI 10.1016/j.neuron.2023.10.021. — URL: \url{https://prescottlab.ca/images/pdf/2023Sagalajev.pdf} (дата обращения: 27.09.2026). — Текст : электронный.
\item[25.] Accelerating with FlyBrainLab the discovery of the functional logic of the Drosophila brain in the connectomic and synaptomic era / A. A. Lazar, T. Liu, M. K. Turkcan, Y. Zhou // eLife. — 2021. — Vol. 10. — Art. e62362. — DOI 10.7554/eLife.62362. — URL: \url{https://elifesciences.org/articles/62362} (дата обращения: 30.09.2026). — Текст : электронный.
\item[26.] AI-based bi-modal fusion system for automated clinical pain monitoring / H. Wang, S. Nienaber, L. Dinges, A. Al-Hamadi // Computers in Biology and Medicine. — 2025. — Vol. 198, no. Pt B. — Art. S0010-4825(25)01614-2. — DOI 10.1016/j.compbiomed.2025.111260. — URL: \url{https://pubmed.ncbi.nlm.nih.gov/41192117/} (дата обращения: 08.10.2026). — Текст : электронный.
\item[27.] AI4PAIN Grand Challenge 2024 dataset : набор данных / R. Fernandez-Rojas, C. Joseph, N. Hirachan [и др.]. — ACIIW 2024 / official challenge. — 2024. — URL: \url{https://sites.google.com/view/ai4pain/challenge-details} (дата обращения: 08.10.2026). — Текст : электронный.
\item[28.] Aktaş, F. A. Explainable AI for pain perception: subject-independent EEG decoding using DeepSHAP and CNNs / F. A. Aktaş, A. Eken, O. Eroğul // Biomedical Physics \& Engineering Express. — 2026. — Vol. 12, no. 1. — Art. 015046. — DOI 10.1088/2057-1976/ae34b4. — URL: \url{https://pubmed.ncbi.nlm.nih.gov/41499809/} (дата обращения: 08.10.2026). — Текст : электронный.
\item[29.] Alhawas, S. Federated Learning in Edge Computing: Vulnerabilities, Attacks, and Defenses—A Survey / S. Alhawas, M. A. Rassam // Sensors. — 2026. — Vol. 26, no. 4. — Art. 1275. — DOI 10.3390/s26041275. — URL: \url{https://www.mdpi.com/1424-8220/26/4/1275} (дата обращения: 25.09.2026). — Текст : электронный.
\item[30.] Aliaga, E. M. R. Emergent Individuality in Whole-Brain Connectome Simulations of Drosophila melanogaster (claude-fly) : программное обеспечение / E. M. R. Aliaga. — GitHub. — 2026. — URL: \url{https://github.com/legacyindiesubmissions-ai/claude-fly} (дата обращения: 08.10.2026). — Текст : электронный.
\item[31.] Aliaga, E. M. R. Emergent Individuality and Neural Integration in Whole-Brain Connectome Simulations of Drosophila melanogaster : препринт / E. M. R. Aliaga. — Zenodo preprint, version 1.0.0 (21 March 2026). — 2026. — DOI 10.5281/zenodo.19152238. — URL: \url{https://zenodo.org/records/19152238} (дата обращения: 24.09.2026). — Текст : электронный.
\item[32.] An Experimental and Clinical Physiological Signal Dataset for Automated Pain Recognition : набор данных / P. Gouverneur, A. Badura, F. Li [и др.]. — Scientific Data. — 2024. — Vol. 11, no. 1. — Art. 1051. — DOI 10.1038/s41597-024-03878-w. — URL: \url{https://figshare.com/articles/dataset/The\_PainMonit\_Database\_An\_Experimental\_and\_Clinical\_Physiological\_Signal\_Dataset\_for\_Automated\_Pain\_Recognition/26965159} (дата обращения: 21.09.2026). — Текст : электронный.
\item[33.] An integrative data-driven model simulating C. elegans brain, body and environment interactions / M. Zhao, N. Wang, X. Jiang [и др.] // Nature Computational Science. — 2024. — Vol. 4, no. 12. — P. 978–990. — DOI 10.1038/s43588-024-00738-w. — URL: \url{https://www.nature.com/articles/s43588-024-00738-w} (дата обращения: 01.10.2026). — Текст : электронный.
\item[34.] An, D. How sparse can a fly-brain model be? Connection strength, wiring placement, network state and task jointly determine function-preserving compression of the Drosophila connectome : препринт / D. An. — bioRxiv. — 2026. — DOI 10.64898/2026.09.19.752860. — URL: \url{https://www.biorxiv.org/content/10.64898/2026.09.19.752860v2.full.pdf} (дата обращения: 29.09.2026). — Текст : электронный.
\item[35.] Anatomical Data Driven Modeling of Evoked Compound Action Potentials Recordings During Spinal Cord Stimulation in a Swine Model / M. K. Brucker-Hahn, A. Deshmukh, M. Settell [и др.] // Neuromodulation: Technology at the Neural Interface. — 2026. — Vol. 29, no. 1. — P. 40–55. — DOI 10.1016/j.neurom.2025.06.008. — URL: \url{https://pubmed.ncbi.nlm.nih.gov/40767809/} (дата обращения: 08.10.2026). — Текст : электронный.
\item[36.] Antropik, M. Fly Brain snnTorch : программное обеспечение / M. Antropik, Neuromorphicism. — GitHub. — 2026. — URL: \url{https://github.com/Neuromorphicism/fly-brain-snntorch} (дата обращения: 08.10.2026). — Текст : электронный.
\item[37.] Artificial Intelligence in Spinal Cord Stimulation and Neuromodulation: A Narrative Review of Clinical Applications, Emerging Evidence, and Future Directions / C. Kolla, S. Madavi, S. Banik [и др.] // JOURNAL OF CLINICAL AND DIAGNOSTIC RESEARCH. — 2026. — Vol. 20, no. 9. — P. UE01. — DOI 10.7860/JCDR/2026/86004.24236. — URL: \url{https://www.jcdr.net/article\_fulltext.asp?id=24236\&issn=0973-709x\&issue=9\&page=UE01\&volume=20\&year=2026} (дата обращения: 28.09.2026). — Текст : электронный.
\item[38.] Artificial Intelligence-Driven Neuromodulation in Neurodegenerative Disease: Precision in Chaos, Learning in Loss / A. Calderone, D. Latella, E. La Fauci [и др.] // Biomedicines. — 2025. — Vol. 13, no. 9. — P. 2118. — DOI 10.3390/biomedicines13092118. — URL: \url{https://www.ebi.ac.uk/europepmc/webservices/rest/PMC12467075/fullTextXML} (дата обращения: 28.09.2026). — Текст : электронный.
\item[39.] Ascending nociceptive pathways drive rapid escape and sustained avoidance in adult Drosophila : препринт / J. M. Jones, A. Sustar, A. Mamiya [и др.]. — bioRxiv. — 2025. — DOI 10.1101/2025.10.28.684868. — URL: \url{https://pubmed.ncbi.nlm.nih.gov/41280033/} (дата обращения: 08.10.2026). — Текст : электронный.
\item[40.] Balboa Binds to Pickpocket In Vivo and Is Required for Mechanical Nociception in Drosophila Larvae / S. E. Mauthner, R. Y. Hwang, A. H. Lewis [и др.] // Current Biology. — 2014. — Vol. 24, no. 24. — P. 2920–2925. — DOI 10.1016/j.cub.2014.10.038. — URL: \url{https://pmc.ncbi.nlm.nih.gov/articles/PMC4438766/} (дата обращения: 08.10.2026). — Текст : электронный.
\item[41.] Baru, M. Controlling neural stimulation using evoked compound action potential signals : патентная заявка WO2026037671A1, международная заявка PCT / M. Baru, F. Meadows. — 2026. — Заявка № PCT/EP2025/072454; заявл. 05.08.2025; опубл. 19.02.2026; заявитель Biotronik SE and Co KG. — URL: \url{https://patents.google.com/patent/WO2026037671A1/en} (дата обращения: 08.10.2026). — Текст : электронный.
\item[42.] Beiran, M. Prediction of neural activity in connectome-constrained recurrent networks / M. Beiran, A. Litwin-Kumar // Nature Neuroscience. — 2025. — Vol. 28, no. 12. — P. 2561–2574. — DOI 10.1038/s41593-025-02080-4. — URL: \url{https://pmc.ncbi.nlm.nih.gov/articles/PMC12648571/} (дата обращения: 24.09.2026). — Текст : электронный.
\item[43.] Betzel, R. F. Diversity of meso-scale architecture in human and non-human connectomes / R. F. Betzel, J. D. Medaglia, D. S. Bassett // Nature Communications. — 2018. — Vol. 9, no. 1. — P. 346. — DOI 10.1038/s41467-017-02681-z. — URL: \url{https://www.nature.com/articles/s41467-017-02681-z} (дата обращения: 08.10.2026). — Текст : электронный.
\item[44.] Biological Processing Units: Leveraging an Insect Connectome to Pioneer Biofidelic Neural Architectures / S. Yu, Z. Qin, T. Liu [и др.] // Lecture Notes in Computer Science. — 2026. — P. 361–369. — DOI 10.1007/978-3-032-00800-8\_32. — arXiv:2507.10951. — URL: \url{https://link.springer.com/chapter/10.1007/978-3-032-00800-8\_32} (дата обращения: 30.09.2026). — Текст : электронный.
\item[45.] BioVid Heat Pain Database : набор данных / S. Walter, S. Gruss, H. Ehleiter [и др.]. — Official database page; Cyberworlds 2013 descriptor. — 2013. — URL: \url{https://www.nit.ovgu.de/BioVid.html} (дата обращения: 08.10.2026). — Текст : электронный.
\item[46.] Biphasic Electrostimulation Artifact Model for ECAP Extraction in Spinal Cord Stimulation / H. Zhang, X. Luo, B. Ma [и др.] // IEEE Transactions on Neural Systems and Rehabilitation Engineering. — 2026. — Vol. 34. — P. 880–893. — DOI 10.1109/TNSRE.2026.3658636. — URL: \url{https://pubmed.ncbi.nlm.nih.gov/41605141/} (дата обращения: 21.09.2026). — Текст : электронный.
\item[47.] Boivin, J. C. Nociception in fruit fly larvae / J. C. Boivin, J. Zhu, T. Ohyama // Frontiers in Pain Research. — 2023. — Vol. 4. — Art. 1076017. — DOI 10.3389/fpain.2023.1076017. — URL: \url{https://www.frontiersin.org/journals/pain-research/articles/10.3389/fpain.2023.1076017/full} (дата обращения: 08.10.2026). — Текст : электронный.
\item[48.] CaMPARI2 enables stimulus-locked whole-brain activity mapping at cellular resolution in unrestrained larval zebrafish / K. R. Robbins, A. Bredbenner, R. A. Osbaldeston [и др.] // Frontiers in Molecular Neuroscience. — 2026. — Vol. 19. — Art. 1772915. — DOI 10.3389/fnmol.2026.1772915. — URL: \url{https://www.frontiersin.org/journals/molecular-neuroscience/articles/10.3389/fnmol.2026.1772915/full} (дата обращения: 01.10.2026). — Текст : электронный.
\item[49.] Cascades and convergence: Dynamic signal flow in a synapse-level brain network / R. F. Betzel, M. G. Puxeddu, C. Seguin, B. Misic // PLOS Complex Systems. — 2026. — Vol. 3, no. 3. — Art. e0000091. — DOI 10.1371/journal.pcsy.0000091. — URL: \url{https://journals.plos.org/complexsystems/article?id=10.1371/journal.pcsy.0000091} (дата обращения: 08.10.2026). — Текст : электронный.
\item[50.] Central versus peripheral neural control of a coordinated walking pattern in Drosophila : препринт / N. Sapkal, D. S. Kumar, S. Sunke [и др.]. — bioRxiv. — 2026. — DOI 10.64898/2026.04.29.721658. — URL: \url{https://www.biorxiv.org/content/10.64898/2026.04.29.721658v2.full} (дата обращения: 24.09.2026). — Текст : электронный.
\item[51.] Centralized brain networks controlling antennal grooming coordination / P. G. Özdil, J. Arreguit, C. Scherrer [и др.] // Nature Communications. — 2026. — Vol. 17, no. 1. — Art. 5617. — DOI 10.1038/s41467-026-72152-x. — URL: \url{https://www.nature.com/articles/s41467-026-72152-x} (дата обращения: 25.09.2026). — Текст : электронный.
\item[52.] Characterization and applications of evoked responses during epidural electrical stimulation / N. Verma, B. Romanauski, D. Lam [и др.] // Bioelectronic Medicine. — 2023. — Vol. 9, no. 1. — Art. 5. — DOI 10.1186/s42234-023-00106-5. — URL: \url{https://link.springer.com/article/10.1186/s42234-023-00106-5} (дата обращения: 25.09.2026). — Текст : электронный.
\item[53.] Choi, K. Systems and methods for providing neurostimulation therapy using multi-dimensional patient features : патентная заявка US20230123383A1, США / K. Choi, E. Ross, D. Whitmer. — 2023. — Заявка № US17/968,734; заявл. 18.10.2022; опубл. 20.04.2023; заявитель Advanced Neuromodulation Systems Inc. — URL: \url{https://patents.google.com/patent/US20230123383A1/en} (дата обращения: 08.10.2026). — Текст : электронный.
\item[54.] Choice of spinal cord stimulation versus targeted drug delivery in the management of chronic pain: a predictive formula for outcomes / N. Mekhail, D. S. Mehanny, S. Armanyous [и др.] // Regional Anesthesia \& Pain Medicine. — 2020. — Vol. 45, no. 2. — P. 131–136. — DOI 10.1136/rapm-2019-100859. — URL: \url{https://pubmed.ncbi.nlm.nih.gov/31932490/} (дата обращения: 24.09.2026). — Текст : электронный.
\item[55.] Classification of chronic pain and spinal cord stimulation response using machine learning in magnetoencephalography data / B. Witjes, M. P. A. Starmans, F. J. P. M. Huygen, C. C. de Vos // PLOS One. — 2025. — Vol. 20, no. 12. — Art. e0337726. — DOI 10.1371/journal.pone.0337726. — URL: \url{https://journals.plos.org/plosone/article?id=10.1371/journal.pone.0337726} (дата обращения: 22.09.2026). — Текст : электронный.
\item[56.] Classification of electrically-evoked potentials in the parkinsonian subthalamic nucleus region / J. Rosing, A. Doyle, A. Brinda [и др.] // Scientific Reports. — 2023. — Vol. 13, no. 1. — Art. 2685. — DOI 10.1038/s41598-023-29439-6. — URL: \url{https://pubmed.ncbi.nlm.nih.gov/36792646/} (дата обращения: 08.10.2026). — Текст : электронный.
\item[57.] Clinical utility of ECAP dosing in a real-world population delivered via EVOKE therapy: the ECAP study / J. E. Pope, G. L. Smith, J. A. Goree [и др.] // Regional Anesthesia \& Pain Medicine. — 2025. — P. rapm–2025–107051. — DOI 10.1136/rapm-2025-107051. — URL: \url{https://rapm.bmj.com/content/early/2025/12/01/rapm-2025-107051} (дата обращения: 22.09.2026). — Текст : электронный.
\item[58.] cnqso. Infinite Sugar : программное обеспечение / cnqso. — GitHub. — 2026. — URL: \url{https://github.com/cnqso/infinite-sugar} (дата обращения: 08.10.2026). — Текст : электронный.
\item[59.] Cold sensing by a glutamate receptor drives avoidance behavior in Drosophila larvae : препринт / L. A. Xu, H. J. Son, X. Bai [и др.]. — bioRxiv. — 2025. — DOI 10.1101/2025.10.05.680544. — URL: \url{https://pmc.ncbi.nlm.nih.gov/articles/PMC12632497/} (дата обращения: 28.09.2026). — Текст : электронный.
\item[60.] Combining brain-wide activity imaging with electron microscopy reveals a distributed brain network for processing threatening somatosensory stimuli : препринт / N. Randel, C. Wang, M. S. Clayton [и др.]. — bioRxiv. — 2025. — DOI 10.1101/2025.09.25.678485. — URL: \url{https://doi.org/10.1101/2025.09.25.678485} (дата обращения: 28.09.2026). — Текст : электронный.
\item[61.] Competitive Disinhibition Mediates Behavioral Choice and Sequences in Drosophila / T. Jovanic, C. M. Schneider-Mizell, M. Shao [и др.] // Cell. — 2016. — Vol. 167, no. 3. — P. 858–870.e19. — DOI 10.1016/j.cell.2016.09.009. — URL: \url{https://zenodo.org/api/records/995631/files/article.pdf/content} (дата обращения: 30.09.2026). — Текст : электронный.
\item[62.] Compound action potentials recorded in the human spinal cord during neurostimulation for pain relief / J. L. Parker, D. M. Karantonis, P. S. Single [и др.] // Pain. — 2012. — Vol. 153, no. 3. — P. 593–601. — DOI 10.1016/j.pain.2011.11.023. — URL: \url{https://pubmed.ncbi.nlm.nih.gov/22188868/} (дата обращения: 25.09.2026). — Текст : электронный.
\item[63.] Computational modeling of paresthesia generated by SCS using a percutaneous lead: a proof-of-concept theoretical model based on an arbitrary somatotopic distribution / T. Le Tutour, K. El Houari, M. Billot [и др.] // Biomedical Physics \& Engineering Express. — 2026. — Vol. 12, no. 2. — Art. 027001. — DOI 10.1088/2057-1976/ae4df4. — URL: \url{https://iopscience.iop.org/article/10.1088/2057-1976/ae4df4} (дата обращения: 30.09.2026). — Текст : электронный.
\item[64.] Connectome simulations identify a central pattern generator circuit for fly walking : препринт / S. M. Pugliese, G. M. Chou, E. T. T. Abe [и др.]. — bioRxiv. — 2025. — DOI 10.1101/2025.09.12.675944. — URL: \url{https://pubmed.ncbi.nlm.nih.gov/42094485/} (дата обращения: 23.09.2026). — Текст : электронный.
\item[65.] Connectome-based biophysical modeling of a figure-ground discrimination circuit : препринт / L. Ziyin, Y. Lu, A. Borst [и др.]. — bioRxiv. — 2026. — DOI 10.64898/2026.09.17.752514. — URL: \url{https://www.biorxiv.org/content/10.64898/2026.09.17.752514v1.full} (дата обращения: 30.09.2026). — Текст : электронный.
\item[66.] Connectome-Based Modelling Reveals Orientation Maps in the Drosophila Optic Lobe : препринт / J. N. Liew, S. Lin, B. Chen [и др.]. — arXiv. — 2026. — DOI 10.48550/arXiv.2609.01330. — arXiv:2609.01330v1. — URL: \url{https://arxiv.org/html/2609.01330v1} (дата обращения: 08.10.2026). — Текст : электронный.
\item[67.] Connectome-based reservoir computing with the conn2res toolbox / L. E. Suárez, A. Mihalik, F. Milisav [и др.] // Nature Communications. — 2024. — Vol. 15, no. 1. — Art. 656. — DOI 10.1038/s41467-024-44900-4. — URL: \url{https://www.nature.com/articles/s41467-024-44900-4} (дата обращения: 30.09.2026). — Текст : электронный.
\item[68.] Connectome-Constrained Latent Variable Models of Whole-Brain Neural Activity / L. Mi, R. Xu, S. Prakhya [и др.] // International Conference on Learning Representations (ICLR). — 2022. — URL: \url{https://openreview.net/forum?id=CJzi3dRlJE-} (дата обращения: 08.10.2026). — Текст : электронный.
\item[69.] Connectome-constrained modeling identifies neurons and synapses that sustain spontaneous activity in Drosophila : препринт / Q. Li, W. Ping, K. Zhang, C. Wang. — bioRxiv. — 2026. — DOI 10.64898/2026.08.21.745055. — URL: \url{https://www.biorxiv.org/content/10.64898/2026.08.21.745055v1.full} (дата обращения: 23.09.2026). — Текст : электронный.
\item[70.] Connectome-constrained networks predict neural activity across the fly visual system / J. K. Lappalainen, F. D. Tschopp, S. Prakhya [и др.] // Nature. — 2024. — Vol. 634, no. 8036. — P. 1132–1140. — DOI 10.1038/s41586-024-07939-3. — URL: \url{https://www.nature.com/articles/s41586-024-07939-3} (дата обращения: 24.09.2026). — Текст : электронный.
\item[71.] Connectome-to-Function: Conditional Generative Latent Representations for Reservoir Computing : препринт / Z. Yu, X. Liu, Y. Jia [и др.]. — arXiv. — 2026. — arXiv:2609.06093. — URL: \url{https://arxiv.org/html/2609.06093} (дата обращения: 08.10.2026). — Текст : электронный.
\item[72.] Connectomic reconstruction of a female Drosophila ventral nerve cord / A. Azevedo, E. Lesser, J. S. Phelps [и др.] // Nature. — 2024. — Vol. 631, no. 8020. — P. 360–368. — DOI 10.1038/s41586-024-07389-x. — URL: \url{https://www.nature.com/articles/s41586-024-07389-x} (дата обращения: 08.10.2026). — Текст : электронный.
\item[73.] Control Signals in Closed-Loop Spinal Cord Stimulation in Patients with Chronic Pain: A Scoping Review / P. Dullur, M. Singhal, B. Hong [и др.] // Neuromodulation: Technology at the Neural Interface. — 2026. — Vol. 29, no. 6. — P. 913. — DOI 10.1016/j.neurom.2025.11.011. — URL: \url{https://www.sciencedirect.com/science/article/pii/S1094715925011705} (дата обращения: 22.09.2026). — Текст : электронный.
\item[74.] Correction: Next‑Generation SCS Programming Platform: Enhancing ECAP Fidelity and Objectivity to Improve Patient Experience / D. J. Parker, A. B. Antony, G. L. Smith [и др.] // Pain and Therapy. — 2026. — Vol. 15, no. 3. — P. 921–921. — DOI 10.1007/s40122-026-00838-7. — URL: \url{https://link.springer.com/article/10.1007/s40122-026-00838-7} (дата обращения: 26.09.2026). — Текст : электронный.
\item[75.] Cross-Domain Multimodal Pain Detection Using a DANN–Transformer Architecture with Supervised Contrastive Warmup and Test-Time Augmentation / S. Bouchfar, Y. Bouchfar, H. E. Malali [и др.] // 2026 International Conference on Intelligent Systems and Digital Applications (ISDA). — 2026. — P. 1–8. — DOI 10.1109/ISDA70544.2026.11606012. — URL: \url{https://doi.org/10.1109/ISDA70544.2026.11606012} (дата обращения: 22.09.2026). — Текст : электронный.
\item[76.] Cross-modal modulation gates nociceptive inputs in Drosophila / G. Pan, R. Li, G. Xu [и др.] // Current Biology. — 2023. — Vol. 33, no. 7. — P. 1372–1380.e4. — DOI 10.1016/j.cub.2023.02.032. — URL: \url{https://pmc.ncbi.nlm.nih.gov/articles/PMC10089977/} (дата обращения: 26.09.2026). — Текст : электронный.
\item[77.] Cross-species connectome comparisons reveal the network attributes of memory capacity and time series prediction : препринт / Y. Yan, G. Wu, H. Yu [и др.]. — bioRxiv. — 2025. — DOI 10.1101/2025.10.29.685101. — URL: \url{https://www.biorxiv.org/content/10.1101/2025.10.29.685101v1.full} (дата обращения: 30.09.2026). — Текст : электронный.
\item[78.] CROSSBRAIN consortium. CROSSBRAIN: Distributed and federated cross-modality actuation through advanced nanomaterials and neuromorphic learning / CROSSBRAIN consortium. — European Commission CORDIS project record. — 2026. — DOI 10.3030/101070908. — URL: \url{https://cordis.europa.eu/project/id/101070908/reporting} (дата обращения: 30.09.2026). — Текст : электронный.
\item[79.] Current state of research and future developments of artificial intelligence in pain diagnosis and treatment / Z. Yang, L. Li, S. Bo [и др.] // Journal of Translational Medicine. — 2026. — Vol. 24, no. 1. — Art. 1177. — DOI 10.1186/s12967-026-08529-9. — URL: \url{https://pubmed.ncbi.nlm.nih.gov/42732075/} (дата обращения: 08.10.2026). — Текст : электронный.
\item[80.] Dahat, C. Multimodal Fusion of Physiological and Psychosocial Parameters for Intelligent Pain Intensity Assessment in Sickle Cell Disease Using Ensemble Learning and Lstm Framework / C. Dahat, G. Soni // 2025 International Conference on Intelligent Communication Networks and Computational Techniques (ICICNCT). — 2025. — P. 1–7. — DOI 10.1109/ICICNCT66124.2025.11232711. — URL: \url{https://doi.org/10.1109/ICICNCT66124.2025.11232711} (дата обращения: 27.09.2026). — Текст : электронный.
\item[81.] Decision-centered wearable biosensors for personalized rehabilitation: integrating multimodal monitoring, artificial intelligence, and closed-loop intervention / F. Jia, X. Li, H. Song [и др.] // Frontiers in Bioengineering and Biotechnology. — 2026. — Vol. 14. — Art. 1948186. — DOI 10.3389/fbioe.2026.1948186. — URL: \url{https://www.frontiersin.org/journals/bioengineering-and-biotechnology/articles/10.3389/fbioe.2026.1948186/full} (дата обращения: 28.09.2026). — Текст : электронный.
\item[82.] Deep Learning and Noninvasive Sensors for Detecting Physiological Dysregulation: A Scoping Review / M. G. Garcés, J. C. Montoya, M. I. P. Martínez [и др.] // Journal of Medical Systems. — 2026. — Vol. 50, no. 1. — Art. 14. — DOI 10.1007/s10916-025-02332-7. — URL: \url{https://pubmed.ncbi.nlm.nih.gov/41615529/} (дата обращения: 26.09.2026). — Текст : электронный.
\item[83.] Deep Multi-Head attention network with Subject-Incremental learning for Cross-Subject generalization in objective pain monitoring using multimodal physiological signals / R. C. Joshi, S. Kumar, S. K. Singh Gautam, M. K. Dutta // Biomedical Signal Processing and Control. — 2026. — Vol. 128. — Art. 111261. — DOI 10.1016/j.bspc.2026.111261. — URL: \url{https://www.sciencedirect.com/science/article/abs/pii/S174680942601815X} (дата обращения: 22.09.2026). — Текст : электронный.
\item[84.] Descending GABAergic pathway links brain sugar-sensing to peripheral nociceptive gating in Drosophila / M. Nakamizo-Dojo, K. Ishii, J. Yoshino [и др.] // Nature Communications. — 2023. — Vol. 14, no. 1. — P. 6515. — DOI 10.1038/s41467-023-42202-9. — URL: \url{https://www.nature.com/articles/s41467-023-42202-9} (дата обращения: 08.10.2026). — Текст : электронный.
\item[85.] Descending neurons integrate learnt information from mushroom body with context to promote escape behaviour : препринт / B. M. W. Jones, S. N. Harris, N. G. Ceffa [и др.]. — bioRxiv. — 2025. — DOI 10.1101/2025.09.30.679458. — URL: \url{https://www.biorxiv.org/content/10.1101/2025.09.30.679458v1.full.pdf} (дата обращения: 24.09.2026). — Текст : электронный.
\item[86.] Devadas, G. R. The Mousing Problem: Why Whole-Brain Emulation Makes AI Suffering Measurable, and What Existing Neuroscience Already Tells Us About It : препринт / G. R. Devadas. — Zenodo preprint, 8 March 2026. — 2026. — DOI 10.5281/zenodo.18911337. — URL: \url{https://zenodo.org/records/18911337} (дата обращения: 25.09.2026). — Текст : электронный.
\item[87.] Development of Machine Learning–Based Models to Predict Treatment Response to Spinal Cord Stimulation / A. Hadanny, T. Harland, O. Khazen [и др.] // Neurosurgery. — 2022. — Vol. 90, no. 5. — P. 523–532. — DOI 10.1227/neu.0000000000001855. — URL: \url{https://telkeslab.com/wp-content/uploads/2023/05/Hadanny-Telkes\_2022\_Development-of-Machine-Learning\%E2\%80\%93Based-Models-to-Predict-Treatment-Response-to-Spinal-Cord-Stimulation.pdf} (дата обращения: 24.09.2026). — Текст : электронный.
\item[88.] Dhiman, N. Topological Sensitivity in Connectome-Constrained Neural Networks : препринт / N. Dhiman. — arXiv. — 2026. — arXiv:2604.04033. — URL: \url{https://arxiv.org/abs/2604.04033} (дата обращения: 08.10.2026). — Текст : электронный.
\item[89.] Dhiman, N. Range-Limited Scale Stability of Dynamic Direction Representations in a Connectome-Constrained Fly Visual Model / N. Dhiman, S. Panwar // Neuroinformatics. — 2026. — Vol. 24, no. 3. — P. 54. — DOI 10.1007/s12021-026-09811-3. — URL: \url{https://link.springer.com/article/10.1007/s12021-026-09811-3} (дата обращения: 08.10.2026). — Текст : электронный.
\item[90.] Dinsmoor, D. A. Spinal cord injury therapy based on evoked compound action potentials : патент US12642969B2, США / D. A. Dinsmoor. — 2026; заявитель Medtronic, Inc. — URL: \url{https://patentsgazette.uspto.gov/week22/OG/html/1547-1/US12642969-20260602.html} (дата обращения: 28.09.2026). — Текст : электронный.
\item[91.] Directional leads applied to spinal cord stimulation: A computational modelling study / Z. Wu, N. Wu, P. Wang [и др.] // PLOS One. — 2026. — Vol. 21, no. 4. — Art. e0345287. — DOI 10.1371/journal.pone.0345287. — URL: \url{https://journals.plos.org/plosone/article?id=10.1371/journal.pone.0345287} (дата обращения: 24.09.2026). — Текст : электронный.
\item[92.] Distributed control circuits across a brain-and-cord connectome / A. S. Bates, J. S. Phelps, M. Kim [и др.] // Nature. — 2026. — Vol. 656, no. 8129. — P. 957–970. — DOI 10.1038/s41586-026-10735-w. — URL: \url{https://pubmed.ncbi.nlm.nih.gov/42259917/} (дата обращения: 08.10.2026). — Текст : электронный.
\item[93.] Divergent Connectivity of Homologous Command-like Neurons Mediates Segment-Specific Touch Responses in Drosophila / S. Takagi, B. T. Cocanougher, S. Niki [и др.] // Neuron. — 2017. — Vol. 96, no. 6. — P. 1373–1387.e6. — DOI 10.1016/j.neuron.2017.10.030. — URL: \url{https://doi.org/10.1016/j.neuron.2017.10.030} (дата обращения: 30.09.2026). — Текст : электронный.
\item[94.] Dobes, M. Drosophila emulations are affectively indeterminate / M. Dobes // Journal of Multiscale Neuroscience. — 2026. — Vol. 5, no. 3. — P. 103–107. — DOI 10.56280/1765455593. — URL: \url{https://doi.org/10.56280/1765455593} (дата обращения: 08.10.2026). — Текст : электронный.
\item[95.] Dopaminergic Modulation of Mushroom Body Output Neurons Mediates Nociception-Induced Escape in Drosophila : препринт / C. L. Yang, C. W. Chen, K. L. Feng [и др.]. — eLife reviewed preprint. — 2026. — DOI 10.7554/eLife.110557.1. — URL: \url{https://elifesciences.org/reviewed-preprints/110557} (дата обращения: 08.10.2026). — Текст : электронный.
\item[96.] Drosophila as a Model to Study the Mechanism of Nociception / J. He, B. Li, S. Han [и др.] // Frontiers in Physiology. — 2022. — Vol. 13. — Art. 854124. — DOI 10.3389/fphys.2022.854124. — URL: \url{https://pmc.ncbi.nlm.nih.gov/articles/PMC8996152/} (дата обращения: 08.10.2026). — Текст : электронный.
\item[97.] Drosophila epidermal cells are intrinsically mechanosensitive and modulate nociceptive behavioral outputs / J. Yoshino, S. S. Mali, C. R. Williams [и др.] // eLife. — 2025. — Vol. 13. — Art. RP95379. — DOI 10.7554/eLife.95379. — URL: \url{https://pubmed.ncbi.nlm.nih.gov/40353351/} (дата обращения: 08.10.2026). — Текст : электронный.
\item[98.] Drosophila mechanical nociceptors preferentially sense localized poking / Z. Liu, M. H. Wu, Q. X. Wang [и др.] // eLife. — 2022. — Vol. 11. — Art. e76574. — DOI 10.7554/eLife.76574. — URL: \url{https://cdn.elifesciences.org/articles/76574/elife-76574-v3.xml} (дата обращения: 27.09.2026). — Текст : электронный.
\item[99.] Drosophila melanogaster foraging regulates a nociceptive-like escape behavior through a developmentally plastic sensory circuit / J. S. Dason, A. Cheung, I. Anreiter [и др.] // Proceedings of the National Academy of Sciences. — 2020. — Vol. 117, no. 38. — P. 23286–23291. — DOI 10.1073/pnas.1820840116. — URL: \url{https://www.pnas.org/doi/10.1073/pnas.1820840116} (дата обращения: 25.09.2026). — Текст : электронный.
\item[100.] Drosophila NOMPC is a mechanotransduction channel subunit for gentle-touch sensation / Z. Yan, W. Zhang, Y. He [и др.] // Nature. — 2013. — Vol. 493, no. 7431. — P. 221–225. — DOI 10.1038/nature11685. — URL: \url{https://pmc.ncbi.nlm.nih.gov/articles/PMC3917554/} (дата обращения: 30.09.2026). — Текст : электронный.
\item[101.] Drosophila pain sensitization and modulation unveiled by a novel pain model and analgesic drugs / W. Jang, M. Oh, E. H. Cho [и др.] // PLOS ONE. — 2023. — Vol. 18, no. 2. — Art. e0281874. — DOI 10.1371/journal.pone.0281874. — URL: \url{https://pmc.ncbi.nlm.nih.gov/articles/PMC9934396/} (дата обращения: 22.09.2026). — Текст : электронный.
\item[102.] Du, J. A Wearable EMG-Driven Closed-Loop TENS Platform for Real-Time, Personalized Pain Modulation / J. Du, S. Luo, P. Shi // Sensors. — 2025. — Vol. 25, no. 16. — Art. 5113. — DOI 10.3390/s25165113. — URL: \url{https://pmc.ncbi.nlm.nih.gov/articles/PMC12390560/} (дата обращения: 22.09.2026). — Текст : электронный.
\item[103.] Duan, S. From Synapses to Dynamics: Obtaining Function from Structure in a Connectome Constrained Model of the Head Direction Circuit : препринт / S. Duan, L. L. Dong, I. Fiete. — bioRxiv. — 2025. — DOI 10.1101/2025.05.26.655406. — URL: \url{https://www.biorxiv.org/content/10.1101/2025.05.26.655406v1.full} (дата обращения: 01.10.2026). — Текст : электронный.
\item[104.] Duarte, R. V. Analysis of Psychological Characteristics Impacting Spinal Cord Stimulation Treatment Outcomes: A Prospective Assessment / R. V. Duarte // Pain Physician. — 2015. — Vol. 3;18, no. 3;5. — P. E369–E377. — DOI 10.36076/ppj.2015/18/E369. — URL: \url{https://www.painphysicianjournal.com/current/pdf?article=MjMyOQ\%3D\%3D\&journal=88} (дата обращения: 24.09.2026). — Текст : электронный.
\item[105.] Durability of Clinical and Quality-of-Life Outcomes of Closed-Loop Spinal Cord Stimulation for Chronic Back and Leg Pain / N. Mekhail, R. M. Levy, T. R. Deer [и др.] // JAMA Neurology. — 2022. — Vol. 79, no. 3. — P. 251. — DOI 10.1001/jamaneurol.2021.4998. — URL: \url{https://jamanetwork.com/journals/jamaneurology/fullarticle/2788004} (дата обращения: 25.09.2026). — Текст : электронный.
\item[106.] Durability of Evoked Compound Action Potential (ECAP)-Controlled, Closed-Loop Spinal Cord Stimulation (SCS) in a Real-World European Chronic Pain Population / H. Nijhuis, J. W. Kallewaard, J. van de Minkelis [и др.] // Pain and Therapy. — 2024. — Vol. 13, no. 5. — P. 1119–1136. — DOI 10.1007/s40122-024-00628-z. — URL: \url{https://link.springer.com/article/10.1007/s40122-024-00628-z} (дата обращения: 25.09.2026). — Текст : электронный.
\item[107.] Early outcomes with a flexible ECAP based closed loop using multiplexed spinal cord stimulation waveforms—single-arm study with in-clinic randomized crossover testing / V. Mohabbati, R. Sullivan, J. Yu [и др.] // Pain Medicine. — 2025. — Vol. 26, no. 11. — P. 773–782. — DOI 10.1093/pm/pnaf058. — URL: \url{https://academic.oup.com/painmedicine/article/26/11/773/8133952} (дата обращения: 27.09.2026). — Текст : электронный.
\item[108.] ECAP-Controlled Closed-Loop Spinal Cord Stimulation for Chronic Nonsurgical Refractory Back Pain / C. W. Hunter, J. S. Raskin, N. A. Mekhail [и др.] // Spine. — 2025. — Vol. 50, no. 23. — P. 1637–1647. — DOI 10.1097/BRS.0000000000005445. — URL: \url{https://pmc.ncbi.nlm.nih.gov/articles/PMC12594148/} (дата обращения: 08.10.2026). — Текст : электронный.
\item[109.] ECAP-controlled closed-loop versus open-loop SCS for the treatment of chronic pain: 36-month results of the EVOKE blinded randomized clinical trial / N. A. Mekhail, R. M. Levy, T. R. Deer [и др.] // Regional Anesthesia \& Pain Medicine. — 2024. — Vol. 49, no. 5. — P. 346–354. — DOI 10.1136/rapm-2023-104751. — URL: \url{https://pmc.ncbi.nlm.nih.gov/articles/PMC11103285/} (дата обращения: 25.09.2026). — Текст : электронный.
\item[110.] Effective Relief of Pain and Associated Symptoms With Closed-Loop Spinal Cord Stimulation System: Preliminary Results of the Avalon Study / M. Russo, M. J. Cousins, C. Brooker [и др.] // Neuromodulation: Technology at the Neural Interface. — 2018. — Vol. 21, no. 1. — P. 38–47. — DOI 10.1111/ner.12684. — URL: \url{https://pubmed.ncbi.nlm.nih.gov/28922517/} (дата обращения: 08.10.2026). — Текст : электронный.
\item[111.] El Othmani, O. Zero-shot multimodal pain estimation via synthetic pain simulation and domain-invariant learning / O. El Othmani, S. Naouali // Frontiers in Artificial Intelligence. — 2026. — Vol. 9. — Art. 1827727. — DOI 10.3389/frai.2026.1827727. — URL: \url{https://www.frontiersin.org/journals/artificial-intelligence/articles/10.3389/frai.2026.1827727/full} (дата обращения: 22.09.2026). — Текст : электронный.
\item[112.] Electroencephalography-Based Pain Detection Using Kernel Spectral Connectivity Network with Preserved Spatio-Frequency Interpretability / S. Buitrago-Osorio, J. Gil-González, A. M. Álvarez-Meza [и др.] // Applied Sciences. — 2025. — Vol. 15, no. 9. — P. 4804. — DOI 10.3390/app15094804. — URL: \url{https://www.mdpi.com/2076-3417/15/9/4804} (дата обращения: 22.09.2026). — Текст : электронный.
\item[113.] Electrophysiologic Characteristics of Epidural Spinal Signals in Preclinical Models of Spinal Cord Stimulation / K. Ladner, E. M. Versantvoort, M. E. G. Thijssen [и др.] // Neuromodulation: Technology at the Neural Interface. — 2026. — Vol. 29, no. 6. — P. 937–949. — DOI 10.1016/j.neurom.2026.01.001. — URL: \url{https://pmc.ncbi.nlm.nih.gov/articles/PMC13519873/} (дата обращения: 30.09.2026). — Текст : электронный.
\item[114.] Eon Systems PBC. Emulation of the Drosophila Fly Brain : программное обеспечение / Eon Systems PBC. — GitHub. — 2026. — URL: \url{https://github.com/eonsystemspbc/fly-brain} (дата обращения: 08.10.2026). — Текст : электронный.
\item[115.] Eon Systems. Eon Systems: mouse-brain emulation programme / Eon Systems. — Eon Systems technical update. — 2026. — URL: \url{https://eon.systems/updates/first-multi-behavior-brain-upload} (дата обращения: 08.10.2026). — Текст : электронный.
\item[116.] Epidural spinal cord recordings (ESRs): sources of neural-appearing artifact in stimulation evoked compound action potentials / A. Deshmukh, M. Settell, K. Cheng [и др.] // Journal of Neural Engineering. — 2025. — Vol. 22, no. 1. — Art. 016050. — DOI 10.1088/1741-2552/ad7f8b. — URL: \url{https://iopscience.iop.org/article/10.1088/1741-2552/ad7f8b} (дата обращения: 25.09.2026). — Текст : электронный.
\item[117.] EPOC: A 28-nm 5.3 pJ/SOP Event-Driven Parallel Neuromorphic Hardware With Neuromodulation-Based Online Learning / F. Chen, Q. Tian, L. Xie [и др.] // IEEE Transactions on Biomedical Circuits and Systems. — 2025. — Vol. 19, no. 3. — P. 629–644. — DOI 10.1109/TBCAS.2024.3470520. — URL: \url{https://pubmed.ncbi.nlm.nih.gov/39356594/} (дата обращения: 24.09.2026). — Текст : электронный.
\item[118.] Error in Figure // JAMA Neurology. — 2022. — Vol. 79, no. 4. — P. 420. — DOI 10.1001/jamaneurol.2022.0022. — URL: \url{https://jamanetwork.com/journals/jamaneurology/fullarticle/2789149} (дата обращения: 26.09.2026). — Текст : электронный.
\item[119.] Estimation and Localization of Chronic Pain level from EEG Signals / S. Mukhopadhyay, M. Sarkar, S. Dey [и др.] // Proceedings of the 24th Annual International Conference on Mobile Systems, Applications and Services Workshops. — 2026. — P. 130–135. — DOI 10.1145/3812836.3814769. — URL: \url{https://doi.org/10.1145/3812836.3814769} (дата обращения: 25.09.2026). — Текст : электронный.
\item[120.] Evoked compound action potential (ECAP)-controlled closed-loop spinal cord stimulation in an experimental model of neuropathic pain in rats / E. M. Versantvoort, B. E. Dietz, D. Mugan [и др.] // Bioelectronic Medicine. — 2024. — Vol. 10, no. 1. — Art. 2. — DOI 10.1186/s42234-023-00134-1. — URL: \url{https://bioelecmed.biomedcentral.com/counter/pdf/10.1186/s42234-023-00134-1.pdf} (дата обращения: 25.09.2026). — Текст : электронный.
\item[121.] Evoked compound action potentials during spinal cord stimulation: effects of posture and pulse width on signal features and neural activation within the spinal cord / M. K. Brucker-Hahn, H. J. Zander, A. J. Will [и др.] // Journal of Neural Engineering. — 2023. — Vol. 20, no. 4. — Art. 046028. — DOI 10.1088/1741-2552/aceca4. — URL: \url{https://iopscience.iop.org/article/10.1088/1741-2552/aceca4/pdf} (дата обращения: 24.09.2026). — Текст : электронный.
\item[122.] Evoked Potentials Recorded From the Spinal Cord During Neurostimulation for Pain: A Computational Modeling Study / C. J. Anaya, H. J. Zander, R. D. Graham [и др.] // Neuromodulation: Technology at the Neural Interface. — 2020. — Vol. 23, no. 1. — P. 64–73. — DOI 10.1111/ner.12965. — URL: \url{https://pmc.ncbi.nlm.nih.gov/articles/PMC6920600/} (дата обращения: 24.09.2026). — Текст : электронный.
\item[123.] Exploiting Large Neuroimaging Datasets to Create Connectome-Constrained Approaches for more Robust, Efficient, and Adaptable Artificial Intelligence : препринт / E. C. Johnson, B. S. Robinson, G. K. Vallabha [и др.]. — arXiv. — 2023. — arXiv:2305.17300. — URL: \url{https://arxiv.org/abs/2305.17300} (дата обращения: 08.10.2026). — Текст : электронный.
\item[124.] External validation of EEG-based machine learning models for continuous pain prediction / T. Mari, J. Henderson, S. H. Ali [и др.] // The Journal of Pain. — 2026. — Vol. 49. — Art. 106416. — DOI 10.1016/j.jpain.2026.106416. — URL: \url{https://pubmed.ncbi.nlm.nih.gov/42600964/} (дата обращения: 08.10.2026). — Текст : электронный.
\item[125.] Factors affecting evoked compound action potential threshold estimation in preclinical models of spinal cord stimulation / E. M. Versantvoort, K. Ladner, D. Mugan [и др.] // Neuromodulation: Technology at the Neural Interface. — 2026. — DOI 10.1016/j.neurom.2026.09.013. — URL: \url{https://doi.org/10.1016/j.neurom.2026.09.013} (дата обращения: 30.09.2026). — Текст : электронный.
\item[126.] Feature extraction and prediction of spinal cord stimulation evoked compound action potentials in humans / S. König, A. Ramadan, D. Sullivan [и др.] // Journal of Neural Engineering. — 2025. — Vol. 22, no. 2. — Art. 026051. — DOI 10.1088/1741-2552/adbfbe. — URL: \url{https://pubmed.ncbi.nlm.nih.gov/40073452/} (дата обращения: 08.10.2026). — Текст : электронный.
\item[127.] Feedback inhibition by a descending GABAergic neuron regulates timing of escape behavior in Drosophila larvae / J. Zhu, J. C. Boivin, A. Garner [и др.] // eLife. — 2024. — Vol. 13. — Art. RP93978. — DOI 10.7554/eLife.93978. — URL: \url{https://www.ebi.ac.uk/europepmc/webservices/rest/PMC11357356/fullTextXML} (дата обращения: 01.10.2026). — Текст : электронный.
\item[128.] Feng, H. A brain-inspired robot pain model based on a spiking neural network / H. Feng, Y. Zeng // Frontiers in Neurorobotics. — 2022. — Vol. 16. — Art. 1025338. — DOI 10.3389/fnbot.2022.1025338. — URL: \url{https://pmc.ncbi.nlm.nih.gov/articles/PMC9807619/} (дата обращения: 22.09.2026). — Текст : электронный.
\item[129.] First evidence of a biomarker-based dose-response relationship in chronic pain using physiological closed-loop spinal cord stimulation / L. Muller, J. Pope, P. Verrills [и др.] // Regional Anesthesia \& Pain Medicine. — 2025. — Vol. 50, no. 4. — P. 345–351. — DOI 10.1136/rapm-2024-105346. — URL: \url{https://pmc.ncbi.nlm.nih.gov/articles/PMC12015077/} (дата обращения: 30.09.2026). — Текст : электронный.
\item[130.] Florea, C. Learning Pain from Emotion: Transferred HoT Data Representation for Pain Intensity Estimation / C. Florea, L. Florea, C. Vertan // Lecture Notes in Computer Science. — 2015. — P. 778–790. — DOI 10.1007/978-3-319-16199-0\_54. — URL: \url{https://link.springer.com/chapter/10.1007/978-3-319-16199-0\_54} (дата обращения: 27.09.2026). — Текст : электронный.
\item[131.] Florea, R. Machine learning for discovery of clinical pain biomarkers following spinal cord injury / R. Florea, K. S. Jeong, C. Y. Saab // Experimental Neurology. — 2026. — Vol. 398. — Art. 115649. — DOI 10.1016/j.expneurol.2026.115649. — URL: \url{https://pubmed.ncbi.nlm.nih.gov/41539462/} (дата обращения: 08.10.2026). — Текст : электронный.
\item[132.] FlyWire FAFB v783 adult female brain connectome : набор данных / S. Dorkenwald, A. Matsliah, A. R. Sterling [и др.]. — FlyWire Codex, версия 783. — 2024. — URL: \url{https://codex.flywire.ai/api/download?data\_version=783} (дата обращения: 08.10.2026). — Текст : электронный.
\item[133.] FlyWire: online community for whole-brain connectomics / S. Dorkenwald, C. E. McKellar, T. Macrina [и др.] // Nature Methods. — 2022. — Vol. 19, no. 1. — P. 119–128. — DOI 10.1038/s41592-021-01330-0. — URL: \url{https://pmc.ncbi.nlm.nih.gov/articles/PMC8903166/} (дата обращения: 21.09.2026). — Текст : электронный.
\item[134.] Fridman, G. Y. Removing artifact in evoked compound action potential recordings in neural stimulators : патент US7835804B2, США / G. Y. Fridman, R. T. Karunasiri. — 2010. — Заявка № US11/734,233; заявл. 11.04.2007; опубл. 16.11.2010; заявитель Advanced Bionics LLC. — URL: \url{https://patents.google.com/patent/US7835804B2/en} (дата обращения: 08.10.2026). — Текст : электронный.
\item[135.] From the fly connectome to exact ring attractor dynamics : препринт / T. Biswas, A. Stanoev, S. Romani, J. E. Fitzgerald. — bioRxiv. — 2024. — DOI 10.1101/2024.11.01.621596. — URL: \url{https://www.biorxiv.org/content/10.1101/2024.11.01.621596v2.full} (дата обращения: 01.10.2026). — Текст : электронный.
\item[136.] gauravvvvvvvvvv. FLYBOX: an open-source connectome sandbox for a simulated fruit-fly nervous system : программное обеспечение / gauravvvvvvvvvv. — GitHub. — 2026. — URL: \url{https://github.com/gauravvvvvvvvvv/flybox} (дата обращения: 08.10.2026). — Текст : электронный.
\item[137.] Generalisation of EEG-Based Pain Biomarker Classification for Predicting Central Neuropathic Pain in Subacute Spinal Cord Injury / K. Anderson, S. Stein, H. Suen [и др.] // Biomedicines. — 2025. — Vol. 13, no. 1. — P. 213. — DOI 10.3390/biomedicines13010213. — URL: \url{https://pmc.ncbi.nlm.nih.gov/articles/PMC11759196/} (дата обращения: 22.09.2026). — Текст : электронный.
\item[138.] GIAFormer: A Gradient-Infused Attention and Transformer for Pain Assessment with EDA-fNIRS Fusion / M. U. Khan, G. Chetty, S. Gkikas [и др.] // Information Fusion. — 2026. — Vol. 131. — Art. 104173. — DOI 10.1016/j.inffus.2026.104173. — URL: \url{https://www.sciencedirect.com/science/article/pii/S1566253526000527} (дата обращения: 22.09.2026). — Текст : электронный.
\item[139.] Gkikas, S. Synthetic Thermal and RGB Videos for Automatic Pain Assessment Utilizing a Vision-MLP Architecture : препринт / S. Gkikas, M. Tsiknakis. — arXiv. — 2024. — P. 4–12. — DOI 10.1109/ACIIW63320.2024.00006. — arXiv:2407.19811. — URL: \url{https://arxiv.org/abs/2407.19811} (дата обращения: 08.10.2026). — Текст : электронный.
\item[140.] Gkikas, S. Efficient Pain Recognition via Respiration Signals: A Single Cross-Attention Transformer Multi-Window Fusion Pipeline : препринт / S. Gkikas, I. Kyprakis, M. Tsiknakis. — arXiv. — 2025. — P. 70–79. — DOI 10.1145/3747327.3764782. — arXiv:2507.21886. — URL: \url{https://arxiv.org/abs/2507.21886} (дата обращения: 25.09.2026). — Текст : электронный.
\item[141.] Gkikas, S. Twins-PainViT: Towards a Modality-Agnostic Vision Transformer Framework for Multimodal Automatic Pain Assessment Using Facial Videos and fNIRS : препринт / S. Gkikas, M. Tsiknakis. — arXiv. — 2024. — P. 13–21. — DOI 10.1109/ACIIW63320.2024.00007. — URL: \url{https://arxiv.org/abs/2407.19809} (дата обращения: 25.09.2026). — Текст : электронный.
\item[142.] Gnana Praveen, R. Deep Weakly Supervised Domain Adaptation for Pain Localization in Videos / R. Gnana Praveen, E. Granger, P. Cardinal // 2020 15th IEEE International Conference on Automatic Face and Gesture Recognition (FG 2020). — 2020. — P. 473–480. — DOI 10.1109/FG47880.2020.00139. — arXiv:1910.08173. — URL: \url{https://doi.org/10.1109/FG47880.2020.00139} (дата обращения: 25.09.2026). — Текст : электронный.
\item[143.] Gualtieri, C. Functional Characterization of Individual Pre- and Postsynaptic Partners In the Drosophila Larval Central Nervous System Using CaMPARI / C. Gualtieri, F. J. Vonhoff // Journal of Visualized Experiments. — 2026, no. 232. — DOI 10.3791/71399. — URL: \url{https://pubmed.ncbi.nlm.nih.gov/42406819/} (дата обращения: 08.10.2026). — Текст : электронный.
\item[144.] Gupta, A. K. Hierarchical Cross-Attention Transformer for Non-contact Multimodal Pain Classification using Remote Physiological Signals and Visual Features / A. K. Gupta, P. Gupta, A. Dhall // 2026 IEEE 20th International Conference on Automatic Face and Gesture Recognition (FG). — 2026. — P. 1–9. — DOI 10.1109/FG67764.2026.11556963. — URL: \url{https://dspace.iiti.ac.in/handle/123456789/18784} (дата обращения: 08.10.2026). — Текст : электронный.
\item[145.] Han, Y. Assessing Pain from Neurophysiological Signals: Machine Learning Approaches Using Functional Connectivity : диссертация / Y. Han. — University of Essex doctoral thesis. — 2024. — URL: \url{https://repository.essex.ac.uk/38532/1/YiyuanHan\_Thesis\_CSEE.pdf} (дата обращения: 24.09.2026). — Текст : электронный.
\item[146.] He, L. The motor pattern of rolling escape locomotion in Drosophila larvae / L. He, L. J. Borjon, W. D. Tracey Jr // Journal of Experimental Biology. — 2026. — Vol. 229, no. 17. — Art. jeb252063. — DOI 10.1242/jeb.252063. — URL: \url{https://www.ebi.ac.uk/europepmc/webservices/rest/PMC13615696/fullTextXML} (дата обращения: 30.09.2026). — Текст : электронный.
\item[147.] Heat-off responses of epidermal cells sensitize Drosophila larvae to noxious inputs : препринт / J. Yoshino, A. Chiu, T. Morita [и др.]. — bioRxiv. — 2025. — DOI 10.1101/2025.05.30.656957. — URL: \url{https://pmc.ncbi.nlm.nih.gov/articles/PMC12157477/} (дата обращения: 28.09.2026). — Текст : электронный.
\item[148.] High-dose spinal cord stimulation for patients with failed back surgery syndrome: a multicenter effectiveness and prediction study / L. Goudman, A. De Smedt, S. Eldabe [и др.] // Pain. — 2021. — Vol. 162, no. 2. — P. 582–590. — DOI 10.1097/j.pain.0000000000002035. — URL: \url{https://pubmed.ncbi.nlm.nih.gov/32910099/} (дата обращения: 08.10.2026). — Текст : электронный.
\item[149.] High-Throughput Analysis of Stimulus-Evoked Behaviors in Drosophila Larva Reveals Multiple Modality-Specific Escape Strategies / T. Ohyama, T. Jovanic, G. Denisov [и др.] // PLoS ONE. — 2013. — Vol. 8, no. 8. — Art. e71706. — DOI 10.1371/journal.pone.0071706. — URL: \url{https://journals.plos.org/plosone/article?id=10.1371/journal.pone.0071706} (дата обращения: 08.10.2026). — Текст : электронный.
\item[150.] HIPPIE: a generative model for electrophysiological analysis across species, technologies, and modalities / J. Gonzalez-Ferrer, J. Lehrer, B. Alvarez-Esteban [и др.] // Nature Communications. — 2026. — Vol. 17, no. 1. — Art. 10007. — DOI 10.1038/s41467-026-76939-w. — URL: \url{https://www.nature.com/articles/s41467-026-76939-w} (дата обращения: 27.09.2026). — Текст : электронный.
\item[151.] Honjo, K. BMP signaling downstream of the Highwire E3 ligase sensitizes nociceptors / K. Honjo, W. D. Tracey // PLOS Genetics. — 2018. — Vol. 14, no. 7. — Art. e1007464. — DOI 10.1371/journal.pgen.1007464. — URL: \url{https://journals.plos.org/plosgenetics/article?id=10.1371/journal.pgen.1007464} (дата обращения: 26.09.2026). — Текст : электронный.
\item[152.] How much of fly walking is written in the wiring? : препринт / I. Guan, Y. Zhao, D. Zhang [и др.]. — arXiv. — 2026. — arXiv:2609.38665. — URL: \url{https://arxiv.org/html/2609.38665v1} (дата обращения: 08.10.2026). — Текст : электронный.
\item[153.] How the Eon Team Produced a Virtual Embodied Fly / S. Harris, A. Sinha, V. Toth [и др.]. — Eon Systems technical update. — 2026. — URL: \url{https://eon.systems/updates/embodied-brain-emulation} (дата обращения: 08.10.2026). — Текст : электронный.
\item[154.] ID\# 1907064 Explainable Machine Learning Pipeline Predicts Spinal Cord Stimulation Responders / I. Telkes, J. Gopal, A. Dunn [и др.] // Neuromodulation: Technology at the Neural Interface. — 2025. — Vol. 28, no. 7. — P. S194. — DOI 10.1016/j.neurom.2025.08.367. — URL: \url{https://www.sciencedirect.com/science/article/abs/pii/S1094715925006439} (дата обращения: 22.09.2026). — Текст : электронный.
\item[155.] ID\# 1907193 Spinal Cord Stimulation (SCS) and Patient Selection: Predictive Modeling in a Global Device Registry / A. Calodney, J. Hatheway, E. Buchser [и др.] // Neuromodulation: Technology at the Neural Interface. — 2025. — Vol. 28, no. 7. — P. S209. — DOI 10.1016/j.neurom.2025.08.396. — URL: \url{https://www.sciencedirect.com/science/article/pii/S1094715925006725} (дата обращения: 26.09.2026). — Текст : электронный.
\item[156.] ID: 318164 Using Neural Networks to Detect Evoked Compound Action Potentials Elicited with Epidural Spinal Cord Stimulation / J. Pilitsis, S. Soghomonyan, A. Will [и др.] // Neuromodulation: Technology at the Neural Interface. — 2024. — Vol. 27, no. 7. — P. S176. — DOI 10.1016/j.neurom.2024.06.321. — URL: \url{https://openalex.org/W4403104801} (дата обращения: 22.09.2026). — Текст : электронный.
\item[157.] Identification and function of thermosensory neurons in Drosophila larvae / L. Liu, O. Yermolaieva, W. A. Johnson [и др.] // Nature Neuroscience. — 2003. — Vol. 6, no. 3. — P. 267–273. — DOI 10.1038/nn1009. — URL: \url{https://www.nature.com/articles/nn1009} (дата обращения: 30.09.2026). — Текст : электронный.
\item[158.] Identifying neural substrates of competitive interactions and sequence transitions during mechanosensory responses in Drosophila / J. B. Masson, F. Laurent, A. Cardona [и др.] // PLOS Genetics. — 2020. — Vol. 16, no. 2. — Art. e1008589. — DOI 10.1371/journal.pgen.1008589. — URL: \url{https://www.ebi.ac.uk/europepmc/webservices/rest/PMC7173939/fullTextXML} (дата обращения: 30.09.2026). — Текст : электронный.
\item[159.] Identifying SCS Trial Responders Immediately After Postoperative Programming with ECAP Dose-Controlled Closed-Loop Therapy / J. E. Pope, A. Antony, E. A. Petersen [и др.] // Pain and Therapy. — 2024. — Vol. 13, no. 5. — P. 1173–1185. — DOI 10.1007/s40122-024-00631-4. — URL: \url{https://link.springer.com/article/10.1007/s40122-024-00631-4} (дата обращения: 08.10.2026). — Текст : электронный.
\item[160.] Impact of Evoked Compound Action Potential (ECAP)-Controlled Closed-Loop Spinal Cord Stimulation in Refractory Lumbar Radiculopathy: A Case Report / A. Soin, S. Kassis, R. J. Witte [и др.] // Cureus. — 2026. — Vol. 18, no. 2. — Art. e102799. — DOI 10.7759/cureus.102799. — URL: \url{https://assets.cureus.com/uploads/case\_report/pdf/460390/20260202-302464-fuovtq.pdf} (дата обращения: 22.09.2026). — Текст : электронный.
\item[161.] Improvements in Therapy Experience With Evoked Compound Action Potential Controlled, Closed-Loop Spinal Cord Stimulation—Primary Outcome of the ECHO-MAC Randomized Clinical Trial / A. Will, M. Fishman, D. Schultz [и др.] // The Journal of Pain. — 2024. — Vol. 25, no. 11. — Art. 104646. — DOI 10.1016/j.jpain.2024.104646. — URL: \url{https://pubmed.ncbi.nlm.nih.gov/39094810/} (дата обращения: 22.09.2026). — Текст : электронный.
\item[162.] In-Clinic and At-Home Patient Experiences with Closed-Loop Spinal Cord Stimulation: 12-Month Outcomes from the Closed-Loop Australia Randomized Crossover Trial / V. Mohabbati, R. Sullivan, J. Yu [и др.] // Pain and Therapy. — 2026. — Vol. 15, no. 5. — P. 1673–1684. — DOI 10.1007/s40122-026-00863-6. — URL: \url{https://link.springer.com/article/10.1007/s40122-026-00863-6} (дата обращения: 01.10.2026). — Текст : электронный.
\item[163.] Januszewski, M. 5 amazing visuals show how the male fruit fly's brain map is advancing neuroscience / M. Januszewski, V. Jain. — Google Research blog. — 2026. — URL: \url{https://blog.google/innovation-and-ai/technology/research/male-fruit-fly-brain-map/} (дата обращения: 08.10.2026). — Текст : электронный.
\item[164.] Jawad, T. Selective peripheral nerve recording using simulated human median nerve activity and convolutional neural networks / T. Jawad, R. G. L. Koh, J. Zariffa // BioMedical Engineering OnLine. — 2023. — Vol. 22, no. 1. — Art. 118. — DOI 10.1186/s12938-023-01181-0. — URL: \url{https://pubmed.ncbi.nlm.nih.gov/38062509/} (дата обращения: 08.10.2026). — Текст : электронный.
\item[165.] Jones, J. M. Beyond Reflex: Nociception, Neural Circuits, and the Transformation of Threat into Memory in Drosophila melanogaster : диссертация / J. M. Jones. — Диссертация, University of Washington. — 2026. — URL: \url{https://digital.lib.washington.edu/researchworks/items/622bc2dc-af95-4861-bc13-a91cab65eac8} (дата обращения: 08.10.2026). — Текст : электронный.
\item[166.] Joseph, C. Real-Time Pain Assessment from Electrodermal Activity Using Deep Learning / C. Joseph, M. Ghahramani, R. F. Rojas // Sensors. — 2026. — Vol. 26, no. 10. — Art. 3020. — DOI 10.3390/s26103020. — URL: \url{https://www.mdpi.com/1424-8220/26/10/3020} (дата обращения: 22.09.2026). — Текст : электронный.
\item[167.] Jung, E. EEG-Based Pain Classification via Sample Selection to Mitigate Subjective Label Bias / E. Jung, S. C. Jun, J. An // IEEE Transactions on Neural Systems and Rehabilitation Engineering. — 2026. — Vol. 34. — P. 2480–2490. — DOI 10.1109/TNSRE.2026.3692232. — URL: \url{https://pubmed.ncbi.nlm.nih.gov/42118625/} (дата обращения: 08.10.2026). — Текст : электронный.
\item[168.] Kakaria, K. S. Ring Attractor Dynamics Emerge from a Spiking Model of the Entire Protocerebral Bridge / K. S. Kakaria, B. L. de Bivort // Frontiers in Behavioral Neuroscience. — 2017. — Vol. 11. — P. 8. — DOI 10.3389/fnbeh.2017.00008. — URL: \url{https://www.frontiersin.org/journals/behavioral-neuroscience/articles/10.3389/fnbeh.2017.00008/full} (дата обращения: 28.09.2026). — Текст : электронный.
\item[169.] Kim, M. J. ROS-mediated activation of Drosophila larval nociceptor neurons by UVC irradiation / M. J. Kim, W. A. Johnson // BMC Neuroscience. — 2014. — Vol. 15, no. 1. — Art. 14. — DOI 10.1186/1471-2202-15-14. — URL: \url{https://pmc.ncbi.nlm.nih.gov/articles/PMC3898224/} (дата обращения: 27.09.2026). — Текст : электронный.
\item[170.] Kopel, E. Pre-registered spectral and certified mixing analysis of the male Drosophila central nervous system connectome : препринт / E. Kopel. — arXiv. — 2026. — DOI 10.48550/arXiv.2609.33054. — arXiv:2609.33054. — URL: \url{https://arxiv.org/html/2609.33054v1} (дата обращения: 08.10.2026). — Текст : электронный.
\item[171.] Labeling of active neural circuits in vivo with designed calcium integrators / B. F. Fosque, Y. Sun, H. Dana [и др.] // Science. — 2015. — Vol. 347, no. 6223. — P. 755–760. — DOI 10.1126/science.1260922. — URL: \url{https://pubmed.ncbi.nlm.nih.gov/25678659/} (дата обращения: 01.10.2026). — Текст : электронный.
\item[172.] Laird, J. H. A model of evoked potentials in spinal cord stimulation / J. H. Laird, J. L. Parker // 2013 35th Annual International Conference of the IEEE Engineering in Medicine and Biology Society (EMBC). — 2013. — P. 6555–6558. — DOI 10.1109/EMBC.2013.6611057. — URL: \url{https://pubmed.ncbi.nlm.nih.gov/24111244/} (дата обращения: 25.09.2026). — Текст : электронный.
\item[173.] Lazar, A. A. NeuroGraphBench: Interacting with Drosophila Connectomes at Scale for Exploring the Functional Logic of Neural Circuits : препринт / A. A. Lazar, S. Shukla, Y. Zhou. — bioRxiv. — 2026. — DOI 10.64898/2026.08.22.746456. — URL: \url{https://www.biorxiv.org/content/10.64898/2026.08.22.746456v1.full} (дата обращения: 27.09.2026). — Текст : электронный.
\item[174.] Light-avoidance-mediating photoreceptors tile the Drosophila larval body wall / Y. Xiang, Q. Yuan, N. Vogt [и др.] // Nature. — 2010. — Vol. 468, no. 7326. — P. 921–926. — DOI 10.1038/nature09576. — URL: \url{https://pmc.ncbi.nlm.nih.gov/articles/PMC3026603/} (дата обращения: 08.10.2026). — Текст : электронный.
\item[175.] Litvak, L. M. Signal classification to detect evoked compound action potential features : патентная заявка WO2025224687A1, международная заявка PCT / L. M. Litvak, J. J. Nedrud. — 2025. — Заявка № PCT/IB2025/054319; заявл. 25.04.2025; опубл. 30.10.2025; заявитель Medtronic Inc. — URL: \url{https://patents.google.com/patent/WO2025224687A1/en} (дата обращения: 08.10.2026). — Текст : электронный.
\item[176.] Long-term safety and efficacy of closed-loop spinal cord stimulation to treat chronic back and leg pain (Evoke): a double-blind, randomised, controlled trial / N. Mekhail, R. M. Levy, T. R. Deer [и др.] // The Lancet Neurology. — 2020. — Vol. 19, no. 2. — P. 123–134. — DOI 10.1016/S1474-4422(19)30414-4. — URL: \url{https://pubmed.ncbi.nlm.nih.gov/31870766/} (дата обращения: 08.10.2026). — Текст : электронный.
\item[177.] Low FoxO expression in Drosophila somatosensory neurons protects dendrite growth under nutrient restriction / A. R. Poe, Y. Xu, C. Zhang [и др.] // eLife. — 2020. — Vol. 9. — Art. e53351. — DOI 10.7554/eLife.53351. — URL: \url{https://pmc.ncbi.nlm.nih.gov/articles/PMC7308081/} (дата обращения: 27.09.2026). — Текст : электронный.
\item[178.] Machine Learning Algorithms Provide Greater Prediction of Response to SCS Than Lead Screening Trial: A Predictive AI-Based Multicenter Study / A. Ounajim, M. Billot, L. Goudman [и др.] // Journal of Clinical Medicine. — 2021. — Vol. 10, no. 20. — P. 4764. — DOI 10.3390/jcm10204764. — URL: \url{https://doi.org/10.3390/jcm10204764} (дата обращения: 24.09.2026). — Текст : электронный.
\item[179.] Machine Learning in Spinal Cord Stimulation for Chronic Pain / V. Hariharan, T. A. Harland, C. Young [и др.] // Operative Neurosurgery. — 2023. — Vol. 25, no. 2. — P. 112–116. — DOI 10.1227/ons.0000000000000774. — URL: \url{https://pubmed.ncbi.nlm.nih.gov/37219574/} (дата обращения: 22.09.2026). — Текст : электронный.
\item[180.] Machine learning predicts spinal cord stimulation surgery outcomes and reveals novel neural markers for chronic pain / J. Gopal, J. Bao, T. Harland [и др.] // Scientific Reports. — 2025. — Vol. 15, no. 1. — Art. 9279. — DOI 10.1038/s41598-025-92111-8. — URL: \url{https://pmc.ncbi.nlm.nih.gov/articles/PMC11920397/} (дата обращения: 08.10.2026). — Текст : электронный.
\item[181.] Machine learning to optimize spinal cord stimulation : патент US11666761B2, США / M. A. Moffitt, N. A. Brill, J. Gu [и др.]. — 2023. — Заявка № US17/027,453; заявл. 21.09.2020; опубл. 06.06.2023; заявитель Boston Scientific Neuromodulation Corp. — URL: \url{https://patents.google.com/patent/US11666761B2/en} (дата обращения: 08.10.2026). — Текст : электронный.
\item[182.] Mauthner, S. E. Optogenetic Stimulation of Nociceptive Escape Behaviors inDrosophilaLarvae : препринт / S. E. Mauthner, W. D. Tracey. — bioRxiv. — 2025. — Vol. 2025, no. 4. — P. pdb.prot108128. — DOI 10.1101/pdb.prot108128. — URL: \url{https://pmc.ncbi.nlm.nih.gov/articles/PMC11787400/} (дата обращения: 22.09.2026). — Текст : электронный.
\item[183.] MDNet: A Lightweight Multidomain 1-D CNN for Embedded Pain Assessment Using EDA Signals / S. Aziz, G. Chetty, R. Goecke, R. F. Rojas // IEEE Internet of Things Journal. — 2026. — Vol. 13, no. 9. — P. 19300–19311. — DOI 10.1109/JIOT.2026.3663683. — URL: \url{https://doi.org/10.1109/JIOT.2026.3663683} (дата обращения: 22.09.2026). — Текст : электронный.
\item[184.] MenstruEase : A Closed‐Loop Smart TENS System for Personalized Management of Primary Dysmenorrhea / S. Chen, S. A. Yar, S. Han [и др.] // Advanced Sustainable Systems. — 2026. — Vol. 10, no. 4. — Art. e70449. — DOI 10.1002/adsu.70449. — URL: \url{https://doi.org/10.1002/adsu.70449} (дата обращения: 22.09.2026). — Текст : электронный.
\item[185.] Method and apparatus for adaptive neural interfaces : патентная заявка WO/2026/146221, международная заявка PCT. — 2026. — URL: \url{https://patentscope.wipo.int/search/en/result.jsf?query=adaptive+neural+interface} (дата обращения: 08.10.2026). — Текст : электронный.
\item[186.] Method of the Year 2025: electron microscopy-based connectomics // Nature Methods. — 2025. — Vol. 22, no. 12. — P. 2463–2464. — DOI 10.1038/s41592-025-02988-6. — URL: \url{https://www.nature.com/articles/s41592-025-02988-6} (дата обращения: 08.10.2026). — Текст : электронный.
\item[187.] Methods and system for recording human physiological signals from implantable leads during spinal cord stimulation / A. Ramadan, S. D. König, M. Zhang [и др.] // Frontiers in Pain Research. — 2023. — Vol. 4. — Art. 1072786. — DOI 10.3389/fpain.2023.1072786. — URL: \url{https://www.frontiersin.org/journals/pain-research/articles/10.3389/fpain.2023.1072786/full} (дата обращения: 08.10.2026). — Текст : электронный.
\item[188.] Mirzakhalili, E. An optimization framework for targeted spinal cord stimulation / E. Mirzakhalili, E. R. Rogers, S. F. Lempka // Journal of Neural Engineering. — 2023. — Vol. 20, no. 5. — Art. 056026. — DOI 10.1088/1741-2552/acf522. — URL: \url{https://doi.org/10.1088/1741-2552/acf522} (дата обращения: 28.09.2026). — Текст : электронный.
\item[189.] ModMix: Data Augmentation for Multimodal Pain Detection / M. Erdal, S. Gruss, S. Walter, F. Schwenker // Lecture Notes in Computer Science. — 2025. — P. 145–155. — DOI 10.1007/978-3-031-88220-3\_11. — URL: \url{https://link.springer.com/chapter/10.1007/978-3-031-88220-3\_11} (дата обращения: 24.09.2026). — Текст : электронный.
\item[190.] Multi-Modal Pain Intensity Assessment Based on Physiological Signals: A Deep Learning Perspective / P. Thiam, H. Hihn, D. A. Braun [и др.] // Frontiers in Physiology. — 2021. — Vol. 12. — Art. 720464. — DOI 10.3389/fphys.2021.720464. — URL: \url{https://www.frontiersin.org/journals/physiology/articles/10.3389/fphys.2021.720464/full} (дата обращения: 28.09.2026). — Текст : электронный.
\item[191.] Multimodal pain assessment with transformers / M. Benavent-Lledo, M. D. Lopez-Valle, D. Ortiz-Perez [и др.] // Neurocomputing. — 2026. — Vol. 664. — Art. 132066. — DOI 10.1016/j.neucom.2025.132066. — URL: \url{https://rua.ua.es/handle/10045/160724} (дата обращения: 08.10.2026). — Текст : электронный.
\item[192.] MuSACo: Multimodal Subject-Specific Selection and Adaptation for Expression Recognition with Co-Training : препринт / M. O. Zeeshan, N. Gillet, A. L. Koerich [и др.]. — arXiv. — 2026. — P. 3606–3616. — DOI 10.1109/WACV61042.2026.00352. — arXiv:2508.12522. — URL: \url{https://arxiv.org/abs/2508.12522} (дата обращения: 08.10.2026). — Текст : электронный.
\item[193.] Network structure governs Drosophila brain functionality / X. Zhang, P. Yang, J. Feng [и др.] // Fundamental Research. — 2026. — Vol. 6, no. 3. — P. 1859–1868. — DOI 10.1016/j.fmre.2025.01.017. — URL: \url{https://www.sciencedirect.com/science/article/pii/S2667325825000792} (дата обращения: 28.09.2026). — Текст : электронный.
\item[194.] Neural circuit mechanisms underlying context-specific halting in Drosophila / N. Sapkal, N. Mancini, D. S. Kumar [и др.] // Nature. — 2024. — Vol. 634, no. 8032. — P. 191–200. — DOI 10.1038/s41586-024-07854-7. — URL: \url{https://www.nature.com/articles/s41586-024-07854-7} (дата обращения: 08.10.2026). — Текст : электронный.
\item[195.] Neural Circuitry that Evokes Escape Behavior upon Activation of Nociceptive Sensory Neurons in Drosophila Larvae / J. Yoshino, R. K. Morikawa, E. Hasegawa, K. Emoto // Current Biology. — 2017. — Vol. 27, no. 16. — P. 2499–2504.e3. — DOI 10.1016/j.cub.2017.06.068. — URL: \url{https://www.cell.com/current-biology/fulltext/S0960-9822(17)30805-9} (дата обращения: 27.09.2026). — Текст : электронный.
\item[196.] Neural circuits for peristaltic wave propagation in crawling Drosophila larvae: analysis and modeling / J. Gjorgjieva, J. Berni, J. F. Evers, S. J. Eglen // Frontiers in Computational Neuroscience. — 2013. — Vol. 7. — P. 24. — DOI 10.3389/fncom.2013.00024. — URL: \url{https://pmc.ncbi.nlm.nih.gov/articles/PMC3616270/} (дата обращения: 30.09.2026). — Текст : электронный.
\item[197.] Neural circuits underlying context-dependent competition between defensive actions in Drosophila larvae / M. Lehman, C. Barré, M. A. Hasan [и др.] // Nature Communications. — 2025. — Vol. 16, no. 1. — P. 1120. — DOI 10.1038/s41467-025-56185-2. — URL: \url{https://www.nature.com/articles/s41467-025-56185-2} (дата обращения: 08.10.2026). — Текст : электронный.
\item[198.] Neural encoding of pain is robust within but unstable between individuals / L. Tiemann, F. S. Bott, E. S. May [и др.] // PLOS Biology. — 2026. — Vol. 24, no. 8. — Art. e3003948. — DOI 10.1371/journal.pbio.3003948. — URL: \url{https://journals.plos.org/plosbiology/article?id=10.1371/journal.pbio.3003948} (дата обращения: 22.09.2026). — Текст : электронный.
\item[199.] Neural signal propagation atlas of Caenorhabditis elegans / F. Randi, A. K. Sharma, S. Dvali, A. M. Leifer // Nature. — 2023. — Vol. 623, no. 7986. — P. 406–414. — DOI 10.1038/s41586-023-06683-4. — URL: \url{https://www.nature.com/articles/s41586-023-06683-4} (дата обращения: 28.09.2026). — Текст : электронный.
\item[200.] Neuroengineering Laboratory, EPFL. FlyGym: Simulating embodied sensorimotor control with NeuroMechFly v2 : программное обеспечение / Neuroengineering Laboratory, EPFL. — GitHub. — 2023. — URL: \url{https://github.com/NeLy-EPFL/flygym} (дата обращения: 08.10.2026). — Текст : электронный.
\item[201.] NeuroMechFly v2: simulating embodied sensorimotor control in adult Drosophila / S. Wang-Chen, V. A. Stimpfling, T. K. C. Lam [и др.] // Nature Methods. — 2024. — Vol. 21, no. 12. — P. 2353–2362. — DOI 10.1038/s41592-024-02497-y. — URL: \url{https://flygym.readthedocs.io/latest/index.html} (дата обращения: 21.09.2026). — Текст : электронный.
\item[202.] Neuromorphic neuromodulation: Towards the next generation of closed-loop neurostimulation / L. F. Herbozo Contreras, N. D. Truong, J. K. Eshraghian [и др.] // PNAS Nexus. — 2024. — Vol. 3, no. 11. — Art. pgae488. — DOI 10.1093/pnasnexus/pgae488. — URL: \url{https://pmc.ncbi.nlm.nih.gov/articles/PMC11565243/} (дата обращения: 22.09.2026). — Текст : электронный.
\item[203.] Neuromorphic Simulation of Drosophila Melanogaster Brain Connectome on Loihi 2 : препринт / F. Wang, B. H. Theilman, F. Rothganger [и др.]. — arXiv. — 2025. — DOI 10.48550/arXiv.2508.16792. — arXiv:2508.16792v1. — URL: \url{https://arxiv.org/html/2508.16792v1} (дата обращения: 08.10.2026). — Текст : электронный.
\item[204.] Neuromorphic Technologies for Neuroengineering: From Adaptive Stimulation to SNN-Based Inference and Deployable Biointerfaces / Z. Sun, A. Mu, F. Hao, H. Wang // Sensors. — 2026. — Vol. 26, no. 10. — Art. 3049. — DOI 10.3390/s26103049. — URL: \url{https://pmc.ncbi.nlm.nih.gov/articles/PMC13210900/} (дата обращения: 26.09.2026). — Текст : электронный.
\item[205.] Neuromuscular basis of Drosophila larval rolling escape behavior / P. C. Cooney, Y. Huang, W. Li [и др.] // Proceedings of the National Academy of Sciences. — 2023. — Vol. 120, no. 51. — Art. e2303641120. — DOI 10.1073/pnas.2303641120. — URL: \url{https://pmc.ncbi.nlm.nih.gov/articles/PMC10743538/} (дата обращения: 30.09.2026). — Текст : электронный.
\item[206.] Neuronal processing of noxious thermal stimuli mediated by dendritic Ca2+ influx in Drosophila somatosensory neurons / S. I. Terada, D. Matsubara, K. Onodera [и др.] // eLife. — 2016. — Vol. 5. — Art. e12959. — DOI 10.7554/eLife.12959. — URL: \url{https://elifesciences.org/articles/12959} (дата обращения: 27.09.2026). — Текст : электронный.
\item[207.] Neurophysiological and Psychophysical Biomarkers for Predicting Spinal Cord Stimulation Treatment Response in Chronic Neuropathic Pain: A Systematic Review / C. A. Odonkor, R. Fagbemigun, F. B. Alhassan, A. Abd-Elsayed // Journal of Personalized Medicine. — 2026. — Vol. 16, no. 10. — P. 486. — DOI 10.3390/jpm16100486. — URL: \url{https://www.mdpi.com/2075-4426/16/10/486} (дата обращения: 30.09.2026). — Текст : электронный.
\item[208.] Next-Generation SCS Programming Platform: Enhancing ECAP Fidelity and Objectivity to Improve Patient Experience / D. J. Parker, A. B. Antony, G. L. Smith [и др.] // Pain and Therapy. — 2026. — Vol. 15, no. 2. — P. 465–480. — DOI 10.1007/s40122-025-00808-5. — URL: \url{https://pmc.ncbi.nlm.nih.gov/articles/PMC13009299/} (дата обращения: 25.09.2026). — Текст : электронный.
\item[209.] Nociception and hypersensitivity involve distinct neurons and molecular transducers in Drosophila / P. Gu, F. Wang, Y. Shang [и др.] // Proceedings of the National Academy of Sciences. — 2022. — Vol. 119, no. 12. — Art. e2113645119. — DOI 10.1073/pnas.2113645119. — URL: \url{https://pmc.ncbi.nlm.nih.gov/articles/PMC8944580/} (дата обращения: 08.10.2026). — Текст : электронный.
\item[210.] Nociceptive interneurons control modular motor pathways to promote escape behavior in Drosophila / A. Burgos, K. Honjo, T. Ohyama [и др.] // eLife. — 2018. — Vol. 7. — Art. e26016. — DOI 10.7554/eLife.26016. — URL: \url{https://elifesciences.org/articles/26016} (дата обращения: 24.09.2026). — Текст : электронный.
\item[211.] Nociceptive Neurons Protect Drosophila Larvae from Parasitoid Wasps / R. Y. Hwang, L. Zhong, Y. Xu [и др.] // Current Biology. — 2007. — Vol. 17, no. 24. — P. 2105–2116. — DOI 10.1016/j.cub.2007.11.029. — URL: \url{https://courses.cit.cornell.edu/bionb491/labresources/week7\%2C\%20Mar\%202-4/Hwangetal2007.pdf} (дата обращения: 25.09.2026). — Текст : электронный.
\item[212.] Nociceptor-Enriched Genes Required for Normal Thermal Nociception / K. Honjo, S. E. Mauthner, Y. Wang [и др.] // Cell Reports. — 2016. — Vol. 16, no. 2. — P. 295–303. — DOI 10.1016/j.celrep.2016.06.003. — URL: \url{https://pmc.ncbi.nlm.nih.gov/articles/PMC5333372/} (дата обращения: 08.10.2026). — Текст : электронный.
\item[213.] Non-random brain connectome wiring enables robust and efficient neural network function under high sparsity : препринт / J. McAllister, C. Houghton, J. Wade, C. O’Donnell. — bioRxiv. — 2026. — DOI 10.64898/2026.03.30.715411. — URL: \url{https://www.biorxiv.org/content/10.64898/2026.03.30.715411v1} (дата обращения: 30.09.2026). — Текст : электронный.
\item[214.] Objective, Same-Day SCS Trials with ECAP-Controlled Closed-Loop Therapy: Depth of Response is Maintained from Trial to 12 months / J. E. Pope, C. M. Vu, J. H. Goree [и др.] // Pain and Therapy. — 2026. — Vol. 15, no. 2. — P. 569–583. — DOI 10.1007/s40122-026-00821-2. — URL: \url{https://pmc.ncbi.nlm.nih.gov/articles/PMC13009343/} (дата обращения: 22.09.2026). — Текст : электронный.
\item[215.] Ohashi, H. Novel behavioral assay of noxious heat responses in Drosophila using decapitated flies / H. Ohashi, T. Sakai // 36th Annual Meeting of the Japanese Society for Comparative Physiology and Biochemistry abstracts, Hikaku seiri seikagaku. — 2014. — Vol. 31, no. 4. — P. 151–164. — URL: \url{https://www.jstage.jst.go.jp/article/hikakuseiriseika/31/4/31\_151/\_pdf/-char/ja} (дата обращения: 28.09.2026). — Текст : электронный.
\item[216.] Othmani, O. E. ZSL\_MPE: code for Zero-shot multimodal pain estimation via synthetic pain simulation and domain-invariant learning : программное обеспечение / O. E. Othmani, S. Naouali. — GitHub. — 2026. — URL: \url{https://github.com/oussama123-ai/ZSL\_MPE} (дата обращения: 08.10.2026). — Текст : электронный.
\item[217.] Ódor, G. Differences in the critical dynamics underlying the human and fruit-fly connectome / G. Ódor, G. Deco, J. Kelling // Physical Review Research. — 2022. — Vol. 4, no. 2. — Art. 023057. — DOI 10.1103/PhysRevResearch.4.023057. — URL: \url{https://journals.aps.org/prresearch/abstract/10.1103/PhysRevResearch.4.023057} (дата обращения: 28.09.2026). — Текст : электронный.
\item[218.] Pain Detection in Biophysiological Signals: Knowledge Transfer from Short-Term to Long-Term Stimuli Based on Distance-Specific Segment Selection / T. B. Ricken, P. Bellmann, S. Walter, F. Schwenker // Computers. — 2023. — Vol. 12, no. 4. — P. 71. — DOI 10.3390/computers12040071. — URL: \url{https://www.mdpi.com/2073-431X/12/4/71} (дата обращения: 27.09.2026). — Текст : электронный.
\item[219.] Pain Detection in Biophysiological Signals: Transfer Learning from Short-Term to Long-Term Stimuli Based on Signal Segmentation / T. B. Ricken, P. Bellmann, S. Walter, F. Schwenker // Lecture Notes in Computer Science. — 2023. — P. 394–404. — DOI 10.1007/978-3-031-37660-3\_28. — URL: \url{https://link.springer.com/chapter/10.1007/978-3-031-37660-3\_28} (дата обращения: 08.10.2026). — Текст : электронный.
\item[220.] Pain in 3D: Controllable Generation of Synthetic Faces for Automated Pain Assessment : препринт / X. L. Lin, S. Mehraban, A. Moturu, B. Taati. — arXiv. — 2025. — arXiv:2509.16727. — URL: \url{https://arxiv.org/html/2509.16727} (дата обращения: 08.10.2026). — Текст : электронный.
\item[221.] Pain intensity classification and evaluation of individual differences in subjects based on hybrid CNN–BiLSTM approach / M. Sun, Y. Liu, D. Geng [и др.] // Biomedical Signal Processing and Control. — 2026. — Vol. 119. — Art. 109815. — DOI 10.1016/j.bspc.2026.109815. — URL: \url{https://www.sciencedirect.com/science/article/pii/S1746809426003691} (дата обращения: 22.09.2026). — Текст : электронный.
\item[222.] Paincontrol: identity-preserving pain expression transfer with generative diffusion models / Y. Zarghami, M. Muzammil, V. Adeli [и др.] // BioMedical Engineering OnLine. — 2026. — Vol. 25, no. 1. — Art. 75. — DOI 10.1186/s12938-026-01561-2. — URL: \url{https://link.springer.com/article/10.1186/s12938-026-01561-2} (дата обращения: 25.09.2026). — Текст : электронный.
\item[223.] PainDiffusion: Learning to Express Pain : препринт / Q. T. Dam, T. Tung Nguyen Nguyen, Y. Endo [и др.]. — arXiv. — 2025. — P. 7253–7260. — DOI 10.1109/IROS60139.2025.11247588. — arXiv:2409.11635. — URL: \url{https://arxiv.org/abs/2409.11635} (дата обращения: 08.10.2026). — Текст : электронный.
\item[224.] painless, a Drosophila Gene Essential for Nociception / W. D. Tracey, R. I. Wilson, G. Laurent, S. Benzer // Cell. — 2003. — Vol. 113, no. 2. — P. 261–273. — DOI 10.1016/S0092-8674(03)00272-1. — URL: \url{https://www.cell.com/cell/fulltext/S0092-8674(03)00272-1} (дата обращения: 25.09.2026). — Текст : электронный.
\item[225.] Parameter-Dependent Effects of Spinal Cord Stimulation on Neural Activation and Evoked Compound Action Potentials : препринт / M. K. Brucker-Hahn, H. J. Zander, D. A. Dinsmoor, S. F. Lempka. — bioRxiv. — 2026. — DOI 10.64898/2026.08.03.742575. — URL: \url{https://www.biorxiv.org/content/10.64898/2026.08.03.742575v1.full} (дата обращения: 30.09.2026). — Текст : электронный.
\item[226.] Patel, A. A. Neural substrates of cold nociception in Drosophila larva / A. A. Patel, A. Cardona, D. N. Cox // eLife. — 2025. — Vol. 12. — Art. RP91582. — DOI 10.7554/eLife.91582. — URL: \url{https://pubmed.ncbi.nlm.nih.gov/40512662/} (дата обращения: 08.10.2026). — Текст : электронный.
\item[227.] Pehlevan, C. Integrative neuromechanics of crawling in D. melanogaster larvae / C. Pehlevan, P. Paoletti, L. Mahadevan // eLife. — 2016. — Vol. 5. — Art. e11031. — DOI 10.7554/eLife.11031. — URL: \url{https://elifesciences.org/articles/11031} (дата обращения: 30.09.2026). — Текст : электронный.
\item[228.] Personalized k-fold Cross-Validation Analysis with Transfer from Phasic to Tonic Pain Recognition on X-ITE Pain Database / Y. Wally, Y. Samaha, Z. Yasser [и др.] // Lecture Notes in Computer Science. — 2021. — P. 788–802. — DOI 10.1007/978-3-030-68780-9\_59. — URL: \url{https://link.springer.com/chapter/10.1007/978-3-030-68780-9\_59} (дата обращения: 27.09.2026). — Текст : электронный.
\item[229.] Physiologic closed loop controlled spinal cord stimulation / T. Kapoor, C. Vu, C. Ung [и др.] // A Modern Atlas for Implantable Devices of the Spine, Brain, and Nerve. — 2026. — P. 283–297. — DOI 10.1016/B978-0-443-36528-7.00011-X. — URL: \url{https://www.sciencedirect.com/book/9780443365287/a-modern-atlas-for-implantable-devices-of-the-spine-brain-and-nerve} (дата обращения: 22.09.2026). — Текст : электронный.
\item[230.] Predicting the Response of High Frequency Spinal Cord Stimulation in Patients with Failed Back Surgery Syndrome: A Retrospective Study with Machine Learning Techniques / L. Goudman, J. P. Van Buyten, A. De Smedt [и др.] // Journal of Clinical Medicine. — 2020. — Vol. 9, no. 12. — P. 4131. — DOI 10.3390/jcm9124131. — URL: \url{https://pmc.ncbi.nlm.nih.gov/articles/PMC7767526/} (дата обращения: 08.10.2026). — Текст : электронный.
\item[231.] Prediction of Response to Spinal Cord Stimulation Using Machine Learning Based on Radiomics and Patient-Reported Outcomes / E. J. Lee, M. L. Edgerton, B. Buccilli [и др.] // Neurosurgery. — 2026. — Vol. 98, no. 4. — P. 895–903. — DOI 10.1227/neu.0000000000003715. — URL: \url{https://pubmed.ncbi.nlm.nih.gov/40879389/} (дата обращения: 08.10.2026). — Текст : электронный.
\item[232.] Preserving privacy in big data research: the role of federated learning in spine surgery / H. Shahzad, C. Veliky, H. Le [и др.] // European Spine Journal. — 2024. — Vol. 33, no. 11. — P. 4076–4081. — DOI 10.1007/s00586-024-08172-2. — URL: \url{https://doi.org/10.1007/s00586-024-08172-2} (дата обращения: 08.10.2026). — Текст : электронный.
\item[233.] Principles of Physiologic Closed-Loop Controllers in Neuromodulation / V. S. Marks, J. J. van Rheede, D. M. Karantonis [и др.] // Neuromodulation: Technology at the Neural Interface. — 2026. — Vol. 29, no. 5. — P. 679–700. — DOI 10.1016/j.neurom.2026.03.010. — URL: \url{https://pubmed.ncbi.nlm.nih.gov/42059841/} (дата обращения: 24.09.2026). — Текст : электронный.
\item[234.] Proof of Concept of Natural Language Processing in Predicting Patient-Reported Outcomes in Spinal Cord Stimulation / M. Kabir, T. Medina, S. R. Daulat [и др.] // Neuromodulation: Technology at the Neural Interface. — 2026. — Art. S1094-7159(26)00624-0. — DOI 10.1016/j.neurom.2026.06.464. — URL: \url{https://pubmed.ncbi.nlm.nih.gov/42524811/} (дата обращения: 08.10.2026). — Текст : электронный.
\item[235.] Pseudo-labeling based adaptations of pain domain classifiers / T. B. Ricken, S. Gruss, S. Walter, F. Schwenker // Frontiers in Pain Research. — 2025. — Vol. 6. — Art. 1562099. — DOI 10.3389/fpain.2025.1562099. — URL: \url{https://www.frontiersin.org/journals/pain-research/articles/10.3389/fpain.2025.1562099/full} (дата обращения: 08.10.2026). — Текст : электронный.
\item[236.] Rajasekhar, G. P. Deep domain adaptation with ordinal regression for pain assessment using weakly-labeled videos / G. P. Rajasekhar, E. Granger, P. Cardinal // Image and Vision Computing. — 2021. — Vol. 110. — Art. 104167. — DOI 10.1016/j.imavis.2021.104167. — arXiv:2008.06392. — URL: \url{https://doi.org/10.1016/j.imavis.2021.104167} (дата обращения: 08.10.2026). — Текст : электронный.
\item[237.] Randomized Modality Mixing with Patchwise RBF Networks for Robust Multimodal Pain Recognition / M. Erdal, S. Gruss, S. Walter, F. Schwenker // Computers. — 2026. — Vol. 15, no. 2. — P. 127. — DOI 10.3390/computers15020127. — URL: \url{https://www.mdpi.com/2073-431X/15/2/127} (дата обращения: 25.09.2026). — Текст : электронный.
\item[238.] Rathee, N. Pain detection from facial expressions using domain adaptation technique / N. Rathee, S. Pahal, P. Sheoran // Pattern Analysis and Applications. — 2022. — Vol. 25, no. 3. — P. 567–574. — DOI 10.1007/s10044-021-01025-4. — URL: \url{https://link.springer.com/article/10.1007/s10044-021-01025-4} (дата обращения: 28.09.2026). — Текст : электронный.
\item[239.] RE-HPBS-IPIC: A Resting EEG- and High-Activation Pain Brain Source-Driven Framework for Inter-Subject Pain Intensity Classification / W. Gao, D. Liu, Q. Wang [и др.] // IEEE Journal of Biomedical and Health Informatics. — 2026. — Vol. 30, no. 6. — P. 5039–5051. — DOI 10.1109/JBHI.2026.3650972. — URL: \url{https://pubmed.ncbi.nlm.nih.gov/41564065/} (дата обращения: 22.09.2026). — Текст : электронный.
\item[240.] Recognition and inhibition of dorsal horn nociceptive signals within a closed-loop system / A. Farajidavar, C. E. Hagains, Y. B. Peng [и др.] // 2010 Annual International Conference of the IEEE Engineering in Medicine and Biology. — 2010. — P. 1535–1538. — DOI 10.1109/IEMBS.2010.5626830. — URL: \url{https://pubmed.ncbi.nlm.nih.gov/21096375/} (дата обращения: 08.10.2026). — Текст : электронный.
\item[241.] Redyno, M. A. M. CLOSING THE LOOP ON CHRONIC PAIN WITH ECAP-CONTROLLED CLOSED-LOOP VS OPEN-LOOP SPINAL CORD STIMULATION: SYSTEMATIC AND LONGITUDINAL POOLED ANALYSIS OF THE EVOKE STUDY / M. A. M. Redyno, A. C. Cresma, R. R. A. Putri // JPHV (Journal of Pain, Vertigo and Headache). — 2026. — Vol. 7, no. 1. — P. 37–44. — DOI 10.21776/ub.jphv.2026.007.01.07. — URL: \url{https://jphv.ub.ac.id/index.php/jphv/article/view/306} (дата обращения: 22.09.2026). — Текст : электронный.
\item[242.] Rethinking Pain Assessment: Subjective Scales, Biomarkers, and Multimodal Integration / A. Liu, L. Zhang, E. Xue [и др.] // Journal of Pain Research. — 2026. — Vol. Volume 19. — P. 1–18. — DOI 10.2147/JPR.S617733. — URL: \url{https://www.dovepress.com/rethinking-pain-assessment-subjective-scales-biomarkers-and-multimodal-peer-reviewed-fulltext-article-JPR} (дата обращения: 22.09.2026). — Текст : электронный.
\item[243.] RNA Sequencing Dataset of Drosophila Nociceptor Translatomic Response to Injury : набор данных / C. M. Hale, K. J. Beauchemin, C. L. Brann [и др.]. — Data. — 2025. — Vol. 10, no. 2. — P. 11. — DOI 10.3390/data10020011. — URL: \url{https://pubmed.ncbi.nlm.nih.gov/40855845/} (дата обращения: 08.10.2026). — Текст : электронный.
\item[244.] Robertson, J. L. Larval Defense against Attack from Parasitoid Wasps Requires Nociceptive Neurons / J. L. Robertson, A. Tsubouchi, W. D. Tracey // PLoS ONE. — 2013. — Vol. 8, no. 10. — Art. e78704. — DOI 10.1371/journal.pone.0078704. — URL: \url{https://journals.plos.org/plosone/article?id=10.1371/journal.pone.0078704} (дата обращения: 26.09.2026). — Текст : электронный.
\item[245.] Saluda Medical Pty Ltd. Novel Treatment Delivery of ECAP-controlled Closed-loop SCS for Chronic Pain : запись клинического исследования NCT04662905 / Saluda Medical Pty Ltd. — ClinicalTrials.gov. — 2020. — URL: \url{https://clinicaltrials.gov/study/NCT04662905} (дата обращения: 08.10.2026). — Текст : электронный.
\item[246.] Schmidgall, S. Biological connectomes as a representation for the architecture of artificial neural networks : препринт / S. Schmidgall, C. Schuman, M. Parsa. — arXiv. — 2022. — DOI 10.48550/arXiv.2209.14406. — arXiv:2209.14406. — URL: \url{https://arxiv.org/abs/2209.14406} (дата обращения: 08.10.2026). — Текст : электронный.
\item[247.] Schwider, T. Mouse to Human Cross-Species Transfer Learning for Electrophysiology-to-Transcriptomics Mapping in Cortical GABAergic Interneurons / T. Schwider, R. Ramezani // Neuroinformatics. — 2026. — Vol. 24, no. 3. — Art. 51. — DOI 10.1007/s12021-026-09806-0. — URL: \url{https://link.springer.com/article/10.1007/s12021-026-09806-0} (дата обращения: 30.09.2026). — Текст : электронный.
\item[248.] Sensory integration and neuromodulatory feedback facilitate Drosophila mechanonociceptive behavior / C. Hu, M. Petersen, N. Hoyer [и др.] // Nature Neuroscience. — 2017. — Vol. 20, no. 8. — P. 1085–1095. — DOI 10.1038/nn.4580. — URL: \url{https://pubmed.ncbi.nlm.nih.gov/28604684/} (дата обращения: 08.10.2026). — Текст : электронный.
\item[249.] Serotonergic Modulation Enables Pathway-Specific Plasticity in a Developing Sensory Circuit in Drosophila / T. Kaneko, A. M. Macara, R. Li [и др.] // Neuron. — 2017. — Vol. 95, no. 3. — P. 623–638.e4. — DOI 10.1016/j.neuron.2017.06.034. — URL: \url{https://pmc.ncbi.nlm.nih.gov/articles/PMC5555049/} (дата обращения: 08.10.2026). — Текст : электронный.
\item[250.] Seung, H. S. Predicting visual function by interpreting a neuronal wiring diagram / H. S. Seung // Nature. — 2024. — Vol. 634, no. 8032. — P. 113–123. — DOI 10.1038/s41586-024-07953-5. — URL: \url{https://www.nature.com/articles/s41586-024-07953-5} (дата обращения: 30.09.2026). — Текст : электронный.
\item[251.] Sexual dimorphism in the complete Drosophila male central nervous system connectome / S. Berg, I. R. Beckett, M. Costa [и др.] // Cell. — 2026. — Vol. 189, no. 18. — P. 5504–5526.e15. — DOI 10.1016/j.cell.2026.08.015. — URL: \url{https://www.cell.com/cell/fulltext/S0092-8674(26)00942-6} (дата обращения: 01.10.2026). — Текст : электронный.
\item[252.] Shibusawa, R. Approximate graph-fibration symmetry in connectome-constrained neural networks: cell-type quotients, equivariance, and behavioural bias : препринт / R. Shibusawa. — SSRN Electronic Journal preprint. — 2026. — DOI 10.2139/ssrn.7537393. — URL: \url{https://doi.org/10.2139/ssrn.7537393} (дата обращения: 30.09.2026). — Текст : электронный.
\item[253.] Simulation insights on the compound action potential in multifascicular nerves / J. J. Tharayil, C. Zinno, F. Agnesi [и др.] // PLOS Computational Biology. — 2025. — Vol. 21, no. 9. — Art. e1013452. — DOI 10.1371/journal.pcbi.1013452. — URL: \url{https://journals.plos.org/ploscompbiol/article?id=10.1371/journal.pcbi.1013452} (дата обращения: 27.09.2026). — Текст : электронный.
\item[254.] Simulation study of spinal cord stimulation evoked compound action potential / Z. Guanghao, Z. Cheng, W. Changzhe, H. Xiaolin // Sheng Wu Yi Xue Gong Cheng Xue Za Zhi = Journal of Biomedical Engineering. — 2021. — Vol. 38, no. 2. — P. 232. — DOI 10.7507/1001-5515.202007016. — URL: \url{https://pmc.ncbi.nlm.nih.gov/articles/PMC9927682/} (дата обращения: 08.10.2026). — Текст : электронный.
\item[255.] Single, P. Cause of Pulse Artefacts Inherent to the Electrodes of Neuromodulation Implants / P. Single, J. Scott // IEEE Transactions on Neural Systems and Rehabilitation Engineering. — 2018. — Vol. 26, no. 10. — P. 2078–2083. — DOI 10.1109/TNSRE.2018.2870169. — URL: \url{https://pubmed.ncbi.nlm.nih.gov/30273153/} (дата обращения: 08.10.2026). — Текст : электронный.
\item[256.] Small conductance Ca2+-activated K+ channels induce the firing pause periods during the activation of Drosophila nociceptive neurons / K. Onodera, S. Baba, A. Murakami [и др.] // eLife. — 2017. — Vol. 6. — Art. e29754. — DOI 10.7554/eLife.29754. — URL: \url{https://elifesciences.org/articles/29754} (дата обращения: 08.10.2026). — Текст : электронный.
\item[257.] Sound response mediated by the TRP channels NOMPC, NANCHUNG, and INACTIVE in chordotonal organs of Drosophila larvae / W. Zhang, Z. Yan, L. Y. Jan, Y. N. Jan // Proceedings of the National Academy of Sciences. — 2013. — Vol. 110, no. 33. — P. 13612–13617. — DOI 10.1073/pnas.1312477110. — URL: \url{https://pmc.ncbi.nlm.nih.gov/articles/PMC3746866/} (дата обращения: 30.09.2026). — Текст : электронный.
\item[258.] Spatiotemporal Distribution of Electrically Evoked Spinal Compound Action Potentials During Spinal Cord Stimulation / J. S. Calvert, R. Darie, S. R. Parker [и др.] // Neuromodulation: Technology at the Neural Interface. — 2023. — Vol. 26, no. 5. — P. 961–974. — DOI 10.1016/j.neurom.2022.03.007. — URL: \url{https://pmc.ncbi.nlm.nih.gov/articles/PMC9643656/} (дата обращения: 25.09.2026). — Текст : электронный.
\item[259.] Strahl, S. Stimulus artifact removal for neuronal recordings : патентная заявка WO2010032132A1, международная заявка PCT / S. Strahl. — 2010. — Заявка № PCT/IB2009/006997; заявл. 15.09.2009; опубл. 25.03.2010; заявитель MED EL Elektromedizinische Geraete GmbH. — URL: \url{https://patents.google.com/patent/WO2010032132A1/en} (дата обращения: 08.10.2026). — Текст : электронный.
\item[260.] Subcallosal Cingulate structural connectivity as biomarker for chronic low back pain : препринт / E. Tsolaki, W. Wei, M. Ward [и др.]. — bioRxiv. — 2024. — DOI 10.1101/2024.09.17.24313765. — URL: \url{https://www.medrxiv.org/content/10.1101/2024.09.17.24313765v1} (дата обращения: 27.09.2026). — Текст : электронный.
\item[261.] Syn Pain: A Synthetic Dataset of Pain and Non-Pain Facial Expressions : препринт / B. Taati, M. Muzammil, Y. Zarghami [и др.]. — arXiv. — 2025. — arXiv:2507.19673. — URL: \url{https://arxiv.org/html/2507.19673} (дата обращения: 08.10.2026). — Текст : электронный.
\item[262.] Synaptic density and relative connectivity conservation maintain circuit stability across development : препринт / I. Fritz, F. Wang, R. Chirif [и др.]. — eLife reviewed preprint. — 2025. — DOI 10.7554/eLife.108643.1. — URL: \url{https://elifesciences.org/reviewed-preprints/108643} (дата обращения: 30.09.2026). — Текст : электронный.
\item[263.] System for patient and adaptive therapy modeling and simulation : патентная заявка WO2026087974A1, международная заявка PCT / L. M. Litvak, S. R. Stanslaski, P. K. Yamaguchi [и др.]. — 2026. — Заявка № PCT/IB2025/059734; заявл. 26.09.2025; опубл. 30.04.2026; заявитель Medtronic Inc. — URL: \url{https://patents.google.com/patent/WO2026087974A1/en} (дата обращения: 08.10.2026). — Текст : электронный.
\item[264.] Systematic Review to Identify Patient-Level Predictors of Treatment Response to Spinal Cord Stimulation for Neuropathic Pain for Studies Published From 2012 to 2024 / A. Kansal, S. Copley, R. V. Duarte [и др.] // Neuromodulation: Technology at the Neural Interface. — 2025. — Vol. 28, no. 6. — P. 890–902. — DOI 10.1016/j.neurom.2025.04.004. — URL: \url{https://www.sciencedirect.com/science/article/pii/S1094715925001485} (дата обращения: 30.09.2026). — Текст : электронный.
\item[265.] Systems and methods for providing neurostimulation therapy according to machine learning operations : патентная заявка US20220323766A1, США / N. R. Hughes, R. Nobles, D. DeShazo, B. Mark. — 2022. — Заявка № US17/566,636; заявл. 30.12.2021; опубл. 13.10.2022; заявитель Advanced Neuromodulation Systems Inc. — URL: \url{https://patents.google.com/patent/US20220323766A1/en} (дата обращения: 08.10.2026). — Текст : электронный.
\item[266.] Temporal Cohorts of Lineage-Related Neurons Perform Analogous Functions in Distinct Sensorimotor Circuits / C. C. Wreden, J. L. Meng, W. Feng [и др.] // Current Biology. — 2017. — Vol. 27, no. 10. — P. 1521–1528.e4. — DOI 10.1016/j.cub.2017.04.024. — URL: \url{https://doi.org/10.1016/j.cub.2017.04.024} (дата обращения: 08.10.2026). — Текст : электронный.
\item[267.] The Choice of Spinal Cord Stimulation Versus Targeted Drug Delivery in the Management of Chronic Pain: Validation of an Outcomes Predictive Formula / N. Mekhail, S. Armanyous, E. Templeton [и др.] // Neuromodulation: Technology at the Neural Interface. — 2023. — Vol. 26, no. 6. — P. 1218–1225. — DOI 10.1016/j.neurom.2023.02.083. — URL: \url{https://pubmed.ncbi.nlm.nih.gov/37061895/} (дата обращения: 24.09.2026). — Текст : электронный.
\item[268.] The connectome of an insect brain / M. Winding, B. D. Pedigo, C. L. Barnes [и др.] // Science. — 2023. — Vol. 379, no. 6636. — Art. eadd9330. — DOI 10.1126/science.add9330. — URL: \url{https://pmc.ncbi.nlm.nih.gov/articles/PMC7614541/} (дата обращения: 01.10.2026). — Текст : электронный.
\item[269.] The digital sphinx: Can a worm brain control a fly body? : препринт / B. W. Brunton, E. T. T. Abe, L. J. Hu, J. C. Tuthill. — bioRxiv. — 2026. — DOI 10.64898/2026.03.20.713233. — URL: \url{https://www.biorxiv.org/content/10.64898/2026.03.20.713233v1} (дата обращения: 30.09.2026). — Текст : электронный.
\item[270.] The Drosophila Connectome as a Computational Reservoir for Time-Series Prediction / L. Costi, A. Hadjiivanov, D. Dold [и др.] // Biomimetics. — 2025. — Vol. 10, no. 5. — P. 341. — DOI 10.3390/biomimetics10050341. — URL: \url{https://www.mdpi.com/2313-7673/10/5/341} (дата обращения: 30.09.2026). — Текст : электронный.
\item[271.] The Drosophila ortholog of vertebrate TRPA1 regulates thermotaxis : препринт / M. Rosenzweig, K. M. Brennan, T. D. Tayler [и др.]. — bioRxiv. — 2005. — Vol. 19, no. 4. — P. 419–424. — DOI 10.1101/gad.1278205. — URL: \url{https://pmc.ncbi.nlm.nih.gov/articles/PMC548941/} (дата обращения: 28.09.2026). — Текст : электронный.
\item[272.] The Drosophila Small Conductance Calcium-Activated Potassium Channel Negatively Regulates Nociception / K. C. E. Walcott, S. E. Mauthner, A. Tsubouchi [и др.] // Cell Reports. — 2018. — Vol. 24, no. 12. — P. 3125–3132.e3. — DOI 10.1016/j.celrep.2018.08.070. — URL: \url{https://pmc.ncbi.nlm.nih.gov/articles/PMC6454897/} (дата обращения: 08.10.2026). — Текст : электронный.
\item[273.] The Effect of Spinal Cord Stimulation Frequency on the Neural Response and Perceived Sensation in Patients With Chronic Pain / G. E. Gmel, R. Santos Escapa, J. L. Parker [и др.] // Frontiers in Neuroscience. — 2021. — Vol. 15. — Art. 625835. — DOI 10.3389/fnins.2021.625835. — URL: \url{https://www.frontiersin.org/journals/neuroscience/articles/10.3389/fnins.2021.625835/full} (дата обращения: 25.09.2026). — Текст : электронный.
\item[274.] The Evoked Compound Action Potential as a Predictor for Perception in Chronic Pain Patients: Tools for Automatic Spinal Cord Stimulator Programming and Control / J. G. Pilitsis, K. V. Chakravarthy, A. J. Will [и др.] // Frontiers in Neuroscience. — 2021. — Vol. 15. — Art. 673998. — DOI 10.3389/fnins.2021.673998. — URL: \url{https://pmc.ncbi.nlm.nih.gov/articles/PMC8320888/} (дата обращения: 25.09.2026). — Текст : электронный.
\item[275.] The fly connectome reveals a path to the effectome / D. A. Pospisil, M. J. Aragon, S. Dorkenwald [и др.] // Nature. — 2024. — Vol. 634, no. 8032. — P. 201–209. — DOI 10.1038/s41586-024-07982-0. — URL: \url{https://pubmed.ncbi.nlm.nih.gov/39358526/} (дата обращения: 08.10.2026). — Текст : электронный.
\item[276.] The Ionotropic Receptors IR21a and IR25a mediate cool sensing in Drosophila / L. Ni, M. Klein, K. V. Svec [и др.] // eLife. — 2016. — Vol. 5. — Art. e13254. — DOI 10.7554/eLife.13254. — URL: \url{https://cdn.elifesciences.org/articles/13254/elife-13254-v2.xml} (дата обращения: 30.09.2026). — Текст : электронный.
\item[277.] The Long-Term Response to High-Dose Spinal Cord Stimulation in Patients With Failed Back Surgery Syndrome After Conversion From Standard Spinal Cord Stimulation: An Effectiveness and Prediction Study / M. De Jaeger, L. Goudman, R. Brouns [и др.] // Neuromodulation: Technology at the Neural Interface. — 2021. — Vol. 24, no. 3. — P. 546–555. — DOI 10.1111/ner.13138. — URL: \url{https://pubmed.ncbi.nlm.nih.gov/32166849/} (дата обращения: 08.10.2026). — Текст : электронный.
\item[278.] The optimization of neuroprosthetic interfaces relying on biophysical and surrogate digital twins / C. Verardo, V. Fossati, L. Toni [и др.] // npj Biomedical Innovations. — 2026. — Vol. 3, no. 1. — Art. 28. — DOI 10.1038/s44385-026-00076-8. — URL: \url{https://www.nature.com/articles/s44385-026-00076-8} (дата обращения: 22.09.2026). — Текст : электронный.
\item[279.] The role of Drosophila Piezo in mechanical nociception / S. E. Kim, B. Coste, A. Chadha [и др.] // Nature. — 2012. — Vol. 483, no. 7388. — P. 209–212. — DOI 10.1038/nature10801. — URL: \url{https://pmc.ncbi.nlm.nih.gov/articles/PMC3297676/} (дата обращения: 08.10.2026). — Текст : электронный.
\item[280.] The TRP Channels Pkd2, NompC, and Trpm Act in Cold-Sensing Neurons to Mediate Unique Aversive Behaviors to Noxious Cold in Drosophila / H. N. Turner, K. Armengol, A. A. Patel [и др.] // Current Biology. — 2016. — Vol. 26, no. 23. — P. 3116–3128. — DOI 10.1016/j.cub.2016.09.038. — URL: \url{https://pmc.ncbi.nlm.nih.gov/articles/PMC5140760/} (дата обращения: 28.09.2026). — Текст : электронный.
\item[281.] Therianos, S. A frozen rate operator from the complete larval connectome: degree and weight govern the gross response, exact wiring governs input routing and mushroom-body modes : препринт / S. Therianos. — arXiv. — 2026. — DOI 10.48550/arXiv.2606.17745. — arXiv:2606.17745. — URL: \url{https://arxiv.org/html/2606.17745v2} (дата обращения: 08.10.2026). — Текст : электронный.
\item[282.] Thermosensory and Nonthermosensory Isoforms of Drosophila melanogaster TRPA1 Reveal Heat-Sensor Domains of a ThermoTRP Channel / L. Zhong, A. Bellemer, H. Yan [и др.] // Cell Reports. — 2012. — Vol. 1, no. 1. — P. 43–55. — DOI 10.1016/j.celrep.2011.11.002. — URL: \url{https://pmc.ncbi.nlm.nih.gov/articles/PMC3278078/?report=reader} (дата обращения: 08.10.2026). — Текст : электронный.
\item[283.] Toward Efficient ECG-Based Pain Intensity Recognition: An End-to-End Neural Network Using Multiple Temporal Feature Compression and Fusion / R. Qiu, K. Xie, S. Xie [и др.] // IEEE Internet of Things Journal. — 2025. — Vol. 12, no. 19. — P. 40764–40778. — DOI 10.1109/JIOT.2025.3590401. — URL: \url{https://doi.org/10.1109/JIOT.2025.3590401} (дата обращения: 24.09.2026). — Текст : электронный.
\item[284.] Toward Objectification of Subjective Chronic Pain Based on Implicit Response in Biosignals / H. S. Seok, S. S. Kim, D. W. Kim, H. Shin // IEEE Transactions on Biomedical Engineering. — 2025. — Vol. 72, no. 1. — P. 337–345. — DOI 10.1109/TBME.2024.3452708. — URL: \url{https://pubmed.ncbi.nlm.nih.gov/39222459/} (дата обращения: 22.09.2026). — Текст : электронный.
\item[285.] Towards Generalizable Learning Models for EEG-Based Identification of Pain Perception : препринт / M. Rezzouk, F. Gagnon, A. Champagne [и др.]. — arXiv. — 2025. — P. 1–6. — DOI 10.1109/MLSP62443.2025.11204206. — URL: \url{https://arxiv.org/abs/2508.11691} (дата обращения: 22.09.2026). — Текст : электронный.
\item[286.] Towards Personalized Edge-AI for Medicine: Efficient Neuromorphic Frameworks for Seizure Detection and Prediction : препринт / L. F. Herbozo Contreras, L. Yu, Z. Huang [и др.]. — bioRxiv. — 2025. — DOI 10.1101/2025.09.22.25336341. — URL: \url{https://www.medrxiv.org/content/10.1101/2025.09.22.25336341v2} (дата обращения: 28.09.2026). — Текст : электронный.
\item[287.] Towards Synthetic Data Generation for Improved Pain Recognition in Videos under Patient Constraints : препринт / J. Nasimzada, J. Kleesiek, K. Herrmann [и др.]. — arXiv. — 2024. — DOI 10.48550/arXiv.2409.16382. — arXiv:2409.16382. — URL: \url{https://arxiv.org/html/2409.16382v1} (дата обращения: 08.10.2026). — Текст : электронный.
\item[288.] Transformer and Attention-Based Models for Automated Pain Assessment: A Systematic Review / S. A. Bukhari, M. S. Hameed, B. Michiels [и др.] // IEEE Access. — 2026. — Vol. 14. — P. 108508–108538. — DOI 10.1109/ACCESS.2026.3713429. — URL: \url{https://vbn.aau.dk/en/publications/transformer-and-attention-based-models-for-automated-pain-assessm/} (дата обращения: 08.10.2026). — Текст : электронный.
\item[289.] Trust-gated synthetic EEG augmentation reduces performance drops when generalizing to new patients / D. Choi, C. Yip, A. Choi, J. Park // npj Digital Medicine. — 2026. — Vol. 9, no. 1. — Art. 634. — DOI 10.1038/s41746-026-02778-0. — URL: \url{https://www.nature.com/articles/s41746-026-02778-0} (дата обращения: 28.09.2026). — Текст : электронный.
\item[290.] Ultrasound-induced wireless implantable stimulator for adaptive pain management // Nature Electronics. — 2025. — Vol. 8, no. 5. — P. 384–385. — DOI 10.1038/s41928-025-01377-3. — URL: \url{https://www.nature.com/articles/s41928-025-01377-3} (дата обращения: 08.10.2026). — Текст : электронный.
\item[291.] Vopson, M. From Connectome to Cognition: Building the First Digital Organism / M. Vopson // IPI Letters. — 2026. — P. N1–N4. — DOI 10.59973/ipil.361. — URL: \url{https://doi.org/10.59973/ipil.361} (дата обращения: 08.10.2026). — Текст : электронный.
\item[292.] Wahezi, S. E. Ecap-filtered neuromodulation waveform matching : патентная заявка US20250099764A1, США / S. E. Wahezi. — 2025. — Заявка № US18/972,074; заявл. 06.12.2024; опубл. 27.03.2025; заявитель Individual. — URL: \url{https://patents.google.com/patent/US20250099764A1/en} (дата обращения: 08.10.2026). — Текст : электронный.
\item[293.] Walker, S. R. Connectomic analysis of taste circuits in Drosophila / S. R. Walker, M. Peña-Garcia, A. V. Devineni // Scientific Reports. — 2025. — Vol. 15, no. 1. — Art. 5278. — DOI 10.1038/s41598-025-89088-9. — URL: \url{https://europepmc.org/articles/PMC11821855} (дата обращения: 26.09.2026). — Текст : электронный.
\item[294.] Wang, B. FLYNN: Robust Neural Network for Robot Navigation using Fly Brain Topology : препринт / B. Wang, J. Chen. — arXiv. — 2026. — DOI 10.48550/arXiv.2607.00025. — arXiv:2607.00025. — URL: \url{https://arxiv.org/abs/2607.00025v2} (дата обращения: 08.10.2026). — Текст : электронный.
\item[295.] Wang, Z. Canine EEG helps human: cross-species and cross-modality epileptic seizure detection via multi-space alignment / Z. Wang, S. Li, D. Wu // National Science Review. — 2025. — Vol. 12, no. 6. — Art. nwaf086. — DOI 10.1093/nsr/nwaf086. — arXiv:2412.17842. — URL: \url{https://academic.oup.com/nsr/article/12/6/nwaf086/8052010} (дата обращения: 25.09.2026). — Текст : электронный.
\item[296.] Whole-Brain Connectomic Graph Model Enables Whole-Body Locomotion Control in Fruit Fly : препринт / Z. Jin, Y. Zhu, C. Zhang, Y. Sui. — arXiv. — 2026. — arXiv:2602.17997. — URL: \url{https://arxiv.org/abs/2602.17997} (дата обращения: 08.10.2026). — Текст : электронный.
\item[297.] Wormuth, A. DoomFly: fly-connectome simulation controlling a Doom arena : программное обеспечение / A. Wormuth. — GitHub. — 2026. — URL: \url{https://github.com/nftechie/doomfly} (дата обращения: 08.10.2026). — Текст : электронный.
\item[298.] Zhang, J. FlyCNS: Connectome-Grounded Information Organization for Communication-Constrained Embodied Control : препринт / J. Zhang, J. Lin, G. Lu. — arXiv. — 2026. — DOI 10.48550/arXiv.2609.28816. — arXiv:2609.28816. — URL: \url{https://arxiv.org/html/2609.28816v1} (дата обращения: 08.10.2026). — Текст : электронный.
\item[299.] Zhong, L. Pickpocket Is a DEG/ENaC Protein Required for Mechanical Nociception in Drosophila Larvae / L. Zhong, R. Y. Hwang, W. D. Tracey // Current Biology. — 2010. — Vol. 20, no. 5. — P. 429–434. — DOI 10.1016/j.cub.2009.12.057. — URL: \url{https://pmc.ncbi.nlm.nih.gov/articles/PMC2995491/} (дата обращения: 25.09.2026). — Текст : электронный.
\item[300.] Zhou, M. G. Representational geometry as a fidelity metric for connectome-constrained networks: evidence from the Drosophila visual system : препринт / M. G. Zhou, J. O. Hasler. — bioRxiv. — 2026. — DOI 10.64898/2026.06.10.731214. — URL: \url{https://www.biorxiv.org/content/10.64898/2026.06.10.731214v11.full} (дата обращения: 30.09.2026). — Текст : электронный.
\end{list}
\endgroup
